% ============================================================
%  Rasmus Nielsen, Grundideernes Logik i kort Begreb.
%  Kjøbenhavn: J. H. Schubothes Boghandel, 1870.  XIV + 89 pp.
%
%  ENGLISH TRANSLATION — from the completed Danish transcription
%  (transcription.tex, the source of truth; image-verified page by page).
%  Title rendered as »The Logic of the Fundamental Ideas in Brief
%  Compass«. Mirrors the transcription's structure 1:1.
%
%  CONVENTIONS (see TRANSLATION-PLAYBOOK.md): moderately literal
%  scholarly register; every \emph{} of the original carried over;
%  Latin kept in \textit{}; Greek glyphs copied verbatim (textalpha);
%  „…“ → ``…''; the id-est mark ɔ: and d.e. → i.e.; em-dash ---.
%  Key terms: Viden = Knowledge; Magt = Power; Sandhed = Truth;
%  Væren/Intet/Vorden = Being/Nothing/Becoming; Forandring = Alteration;
%  Væsen/Skin = Essence/Semblance; Grund = Ground; Vexelvirkning =
%  Reciprocity; det Almene / det Særskilte / det Enkelte = the Universal /
%  the Particular / the Singular; Dom = Judgment; Slutning = Inference
%  (uniformly, incl. Hegel's Schluss contexts); Grunddeling =
%  ground-division; Formaal/Midler = End/Means; Idee = Idea.
%  The Danish printer's errors (»meu«, »tænkcs«, »Saudhed« etc.) are
%  diplomatic facts of the transcription; the translation renders the
%  corrected sense and does not reproduce them.
%
%  STATUS: front matter (Foreword + Introduction, fol. III–XIII)
%  translated; body pp. 1–89 and Contents outstanding.
% ============================================================
\documentclass[12pt]{book}
\usepackage[utf8]{inputenc}
\usepackage[T1]{fontenc}
\usepackage[english]{babel}
\usepackage[protrusion=true,expansion=true]{microtype}
\usepackage[margin=1.4in]{geometry}
\usepackage{setspace}
\usepackage{enumitem}
\usepackage{libertinus}
\usepackage{libertinust1math}
\usepackage{textalpha}   % Greek glyphs copied from the Danish
\usepackage{fancyhdr}
\usepackage[colorlinks=true,linkcolor=black,urlcolor=black,
  pdftitle={The Logic of the Fundamental Ideas in Brief Compass},
  pdfauthor={Rasmus Nielsen}]{hyperref}
\setlength{\headheight}{15pt}
\pagestyle{fancy}
\fancyhf{}
\fancyhead[LE,RO]{\thepage}
\fancyhead[RE]{\textit{The Logic of the Fundamental Ideas}}
\fancyhead[LO]{\textit{\leftmark}}
\setlength{\parindent}{1.5em}

\setstretch{1.3}

\begin{document}

\begin{titlepage}
\centering
{\Huge The Logic of the Fundamental Ideas\par}
\vspace{1ex}
{\Large in Brief Compass.\par}
\vspace{4ex}
{\large By\par}
\vspace{2ex}
{\LARGE R. Nielsen.\par}
\vspace{4ex}
{\large Translated from the Danish\par}
{\large (Kjøbenhavn: J. H. Schubothes Boghandel, 1870)\par}
\end{titlepage}

\frontmatter

\chapter*{Foreword}
\markboth{Foreword}{Foreword}
% --- fol. III ---
In this little book the author has endeavoured to give a short and
clear presentation of his theistic logic. He has therein found
occasion to bring out the nodal points on which the scientific
dispute between pantheism and theism must essentially turn.

\begin{flushright}R. Nielsen.\end{flushright}

% (fol. IV is blank.)

\chapter*{Introduction}
\markboth{Introduction}{Introduction}

% --- fol. V ---
At the present stage of philosophy's development it becomes a
scientific necessity to treat logic from three essentially different
points of view.

In order to develop the laws of thinking, one must keep to thinking
itself. But thinking is a human activity of spirit; we must learn
what thinking is from our own thinking. It lies near, then, to regard
the human being as the thinking subject, to reflect upon the mutual
relations of the acts of thought, and to mark the forms necessary for
the understanding, the intellectual traces which advancing thought
deposits. This point of view is characteristic of \emph{formal
logic}.

It is, however, so far from being the case that the finite
understanding should command the demands of thought, that these must
on the contrary be said to command the thinker; there shall be
thinking, but when the thinking is consistent it is for the rest
indifferent who it is that thinks --- whether it be a human being or
a god or perhaps pure being. The emphasis falls on the thinking. For
thinking's own sake, on the other hand, it is of importance to fix
the eye on that which is thought, on the universally necessary
content; for the content of thought and the form of thought are
determined by one another. This point of view is characteristic of
the \emph{speculative} or, to use a perhaps still more telling
expression, \emph{the immanent logic}.

These two principal forms of logic, of which the first has an
Aristotle, the second a Hegel, as supreme representative, may with
reason be said to enjoy a well-deserved recognition in science.

That they cannot well be reconciled with one another, but stand in a
% --- fol. VI ---
most serious dispute concerning the very first principles of
thinking, notably the principle of contradiction, is certainly less
commendable; but they will both maintain their place: the one will
not be able to displace the other.

In order that each of the contending parties may come fully into its
right, it is in our judgment necessary that yet a third principal
form be introduced. But before we attempt to prove the correctness of
this assertion, we consider it expedient to bring the architectonic
opposition between the two doctrinal edifices nearer to intuition by
prefacing an outline corresponding to each of them.

\section*{Formal Logic.}

Formal logic, or the doctrine of the abstractly necessary
thought-forms, divides into the doctrine of \emph{Concept},
\emph{Judgment} and \emph{Inference}.

\subsection*{Concept.}

Thinking, which is already present in sensation, raises itself forth
as it dissolves the immediate sense-images, forms by comparison and
abstraction a \emph{general image}, holds it fast in the \emph{word},
makes clear to itself the essential marks that necessarily belong
together, and transforms it into a \emph{concept} (concept-analysis,
definition).

The sum of the determinations belonging to a concept constitutes the
concept's \emph{content}; the number of the objects to which it finds
application designates its \emph{extension}. (Inverse relation
between content and extension; superordinate and subordinate
concepts; genus- and species-concepts.)

\subsection*{Judgment.}

The concept finds its pure logical expression only in language: words
are joined into sentences, and in sentences judgments are uttered.
The sentence whose form expresses a concept is a \emph{logical
judgment}. In order to express the concept in this form, the judgment
must formally indicate: the \emph{Subject}, the representation that
is to be determined; the \emph{Predicate}, the determining
representation; and the \emph{Copula}, the determining word of
connection.

From the subject's connection with the predicate, determined by the
copula, it follows that the judgment can be separated into a
plurality of logically different judgments, and judgments are
divided: according to \emph{quality} into affirmative, negative and
infinite; according to \emph{quantity} into singular, particular and
general; according to \emph{relation} into categorical, hypothetical
and disjunctive; according to \emph{modality} into assertoric,
% --- fol. VII ---
problematic and apodictic.

If the predicate proceeds from the subject with universal logical
necessity, the judgment is \emph{analytic}; if the predicate can be
joined to the subject only through a particular intuition or
experience, the judgment is \emph{synthetic} (synthetic judgments
\textit{a priori}; Kant's problem).

\subsection*{Inference.}

The judgment must be grounded; its validity must be seen into; the
act of thinking whereby the validity of one judgment is derived from
one or more other judgments is a \emph{logical inference}. If the
inference has only one antecedent proposition (premise), it is called
\emph{immediate}; if it has several premises, it is called
\emph{mediate} (the requisites of inference).

The mediate inferences are divided into: the \emph{categorical},
conditioned by the subsumption of the subordinate concept's mark
(\textit{nota notæ}) under the superordinate (syllogistic figures);
the \emph{hypothetical}, which rest on the dependence of the
conditioned upon its condition (\textit{modus ponens} and
\textit{modus tollens}); and the \emph{disjunctive}, which are
grounded on this, that the members of a whole are mutually so related
that when one is expressly posited the rest are excluded
(\textit{modus ponendo tollens}), and that when the rest are excluded
the one remaining must be posited (\textit{modus tollendo ponens}).

By treating the doctrine of proof and method, formal logic resolves
itself into a general doctrine of knowledge and of science.

\section*{The Immanent Logic.}

The immanent logic, which is a dialectical development of the system
of the categories (the thought-forms inherent in all being), divides
into the doctrine of \emph{Being}, \emph{Essence} and \emph{Concept}.

\subsection*{Being.}

Universal being is one with universal thinking; but the unity of
thinking and being is, in its most abstract form, void of content:
pure \emph{Being} is the same as pure \emph{Nothing}. But Nothing, as
non-being, is nevertheless different from Being. That Being both is
and is not Nothing means that Being unceasingly passes over into
Nothing, and Nothing into Being: this transition is \emph{Becoming}.

In Becoming, Nothing is a limit for Being, and Being for Nothing.
Being limited by its Nothing is a \emph{being}, a \emph{Something},
% --- fol. VIII ---
and Nothing determined by Being, an \emph{Other}. But the Other is
itself Something, and Something in turn Other; Something passes by
its own limit over into Other, and Other into Something: this
transition is \emph{Alteration}.

In alteration Something is one with its Other; this is the
\emph{Unchangeable} in the One; but Something is also different from
the Other; this is the \emph{Changeable} in the One. It is thus the
Unchangeable that is altered; but precisely by being altered it keeps
% (the Danish prints »meu« for »men« here; see the transcription)
itself unaltered: under all alteration there is something
unchangeable.

\subsection*{Essence.}

The Unchangeable that is seen in the Changeable is \emph{Essence};
the Changeable through which the Essence is seen is \emph{Semblance}:
the alternating shine of Semblance and Essence is \emph{Reflection}.

In its identity with the Essence, the Semblance is
\emph{Phenomenon}; in its difference from the Phenomenon, the Essence
is \emph{Ground}. The Phenomenon goes to the ground; it belongs to
Semblance and must be reflected into its Essence. But in that the
Phenomenon receives the Essence for its ground, it becomes grounded
and is posited as \emph{Existence}.

Existence reflected into itself is an \emph{existent}, a
\emph{Thing}; but the Thing can be reflected into itself only by
reflecting its existence into other existing things: the
determinations corresponding to these reflections are
\emph{Properties}.

In this mutual reflecting the things dissolve into the properties,
and everything becomes unstable. But the unstable presupposes
something subsisting which cannot be dissolved, because its
subsistence infinitely confirms itself in the alternation of the
unstable; this subsisting something is the \emph{Substance}, and the
modes of appearance alternating in instability are the
\emph{Accidents}.

The relation between the Substance and its Accidents is a
\emph{relation of causality}: the Substance is \emph{Cause}, the
original matter; the Accidents are \emph{Effects}. But the Cause has
its substantial expression in the Effect; without Effect, no Cause:
by making the Cause a Cause, the Effect is itself Substance --- thus
a Cause that makes the Cause an Effect.

In \emph{Reciprocity} the Effect is Cause, and the Cause Effect; and
in that all the moments of actuality contained in the Essence are,
under all-sided reciprocity, developed in full totality, the Essence
developed in the totality is Concept.
% --- fol. IX ---

\subsection*{Concept.}

The self-comprehending totality is, in its unity with itself, the
\emph{Universal}; in its difference from itself, \emph{the
Particular}; in its totalized determinateness through itself, the
\emph{Singular}.

The Concept's self-development is its \emph{ground-division} (the
Judgment). By virtue of the ground-division, the totality in its
first immediate being, the Singular, is posited as the
\emph{Subject}, and the totality in the reflected form, the
Universal, as the \emph{Predicate}: the Subject is what the Predicate
is; the \emph{Singular} is the \emph{Universal}.

But in that the Singular is comprehended as the Universal, and the
Universal as the Singular, the opposed members of the
concept-separation become mutually comprehended, in reciprocal
determination, as the \emph{Particular}.

The outer members of the ground-division, the \emph{Extremes}, have
thus in the Concept's self-development deposited a \emph{middle
member}, a \emph{Medium}, in which they can be mediated. This
mediation is the \emph{Inference}.

The Concept's self-mediation through Judgment and Inference is its
\emph{Subjectivity}; the immediacy regained through the mediation is
its \emph{Objectivity}.

As immediate, existing totality the Concept is Object --- and, by
virtue of the ground-division, an infinite manifold of existing
objects.

The objects in their outward, indifferent mutuality are
\emph{mechanical}, acting upon one another by \emph{pressure} and
\emph{impact}. The separation of motion and rest which herewith
enters is a relative appearing of the \emph{absolute Mechanics}, in
which motion is rest, and rest motion.

Beneath the indifference of the mechanical objects there hides a
\emph{difference}, which comes to light in chemical attraction. The
difference of the objects brings with it a tension, which is
equalized in \emph{the chemical product}. But the product, in which
the difference is neutralized, must as the \emph{Neutral} dissolve
again into difference, and the \emph{process} begins anew.

In the chemical process the Concept goes forth out of its own
presuppositions, sublates the presuppositions in their result, in
order again to sublate the result in the presuppositions. The
objective Concept is now master over the material: the result
anticipated in the presuppositions is \emph{End}, the presuppositions
are \emph{Means}, the movement is \emph{teleological}.

The teleological Concept, comprehending itself in its own elements,
is a form that gives itself content, a content that gives itself
form: in this its \emph{absolute unity} of form and content, of
\emph{Subjective and Objective}, of end and means, the Concept is
\emph{Idea}.
% --- fol. X ---

From the outlines here set forth it will be seen how divergent the
directions are. In formal logic the human being, the thinking
subject, is from first to last occupied with his own acts of thought;
it is the human being himself who abstracts, reflects, judges and
infers. But for just that reason the thought-forms are empty forms,
and the acts of thought finite, psychologically conditioned
operations of the understanding. The object of thinking falls outside
the thinking. Here is an insurmountable gulf between the subjective
and the objective, and the thinker has only to watch that in his
judgment about the things he does not come into contradiction with
himself. To secure the unity of the thinkable and the possible, one
holds fast to \emph{the principle of contradiction}; but that the
principle of contradiction affords no reliable guarantee of the
objective validity of subjective thinking, the \emph{antinomies} in
Kant's Critique of Pure Reason have sufficiently shown.

The gulf between the subjective and the objective disappears in the
immanent logic, which abstracts from the thinking subjectivity and
lets Being think itself by the help of contradiction and negation.
That Being is Nothing is indeed a contradiction, for Being must
necessarily exclude non-being. But the non-being that is excluded is
precisely the negation that lies in Being's own concept. That the
Unchangeable is changeable, that the Essence is Semblance, that the
Cause is Effect, are sheer contradictions. But these contradictions
make it precisely evident how the thinking that is entangled in
finite representations deceives itself by keeping the content of what
is outside itself. The Changeable is the Unchangeable's indwelling
negation, just as the Semblance is the Essence's, and the Effect the
Cause's.

Nevertheless it must not be overlooked that the immanent method of
developing the matter, despite its epoch-making significance, also
has its weak sides. Negation is originally a thought-determination,
not a determination of being. What is can, in its immediacy, negate
neither itself nor another being. Negations, contradictions and
resolutions of contradictions are acts of thought which hold of what
is only under the express condition that Being is thought. But in
order that this condition can be fulfilled, the invoked unity of
thinking and being must from the beginning be grounded upon a
corresponding opposition. Pure Being is Nothing, when Being is
thought; but Nothing and Being fall nonetheless each under its own
sphere: Nothing under thinking, and so under the subjective; and
Being under actuality, and so under the objective.

If the dispute between the two recognized points of view is not to
% --- fol. XI ---
become endless and unfruitful, a third point of view must be chosen,
from which it becomes possible to reconcile the contending parties.
It must then be conceded to formal logic that the thoughts proceed
from a thinking subjectivity; but it must at the same time be
conceded to the immanent logic that it is not a finite but an
infinite thinking that comprehends existence. Formal logic must be
right in this, that there is a characteristic difference between the
subjective in thinking and the objective in actuality; but the logic
of immanence retains its right in this, that the opposition of the
subjective and the objective rests upon identity of thinking and
being. It is this third point of view that is chosen in \emph{the
Logic of the Fundamental Ideas}.

As regards the origin of thoughts, only three assumptions are
possible: 1)~that they proceed from a finite subjectivity; 2)~that
they develop themselves out of universal being; 3)~that they spring
from an infinite, absolute subjectivity. As regards the relation
between thinking and being, there are likewise only three assumptions
possible: 1)~that it is a finite opposition, be it in agreement or in
discord; 2)~that the unity sublates the opposition; 3)~that the unity
grounds the opposition.

For the all-sided development of logic the three points of view here
set up are necessary; but more than these three are, so far as
principles are concerned, impossible.

\bigskip
\centerline{\rule{6em}{0.4pt}}
\bigskip

% --- fol. XII ---
\section*{The Logic of the Fundamental Ideas}
% (as in the Danish, the heading opens the sentence that follows)
\noindent develops the ideal essentiality of the thought-forms under
\emph{the Idea of Knowledge}, the forms of being corresponding to
actuality under \emph{the Idea of Power}, and points out the
dualistic unity of the ideal and the real content in \emph{the Idea
of Truth}.

The fundamental separation of the three Ideas is especially motivated
by regard to the one-sidedness of the immanent logic.

In order that thinking shall be possible, there must be something
that is thought, hence something that thinking itself posits,
% (the Danish prints »tænkcs« for »tænkes« here; see the transcription)
presupposes and produces; this is \emph{the Ideal}. But in order that
actuality, the actual world, can have validity, there must be
something other than thinking: something that is not thought, but
about which there is thinking; something whereby thinking is
determined, but which it does not itself produce; this is \emph{the
Real}. The difference between the Ideal and the Real points to a
decisive difference between Knowledge and Power: the Ideal is a
knowledge-determination, the Real a power-determination.

It is universally acknowledged that everything in the world must have
a ground; herein lies that the existent must rest upon a
non-existent. The ground by itself is not existent; it has existence
only in the existence which it grounds. The ground is an
\emph{ideality}; the thing, the actual object, is a \emph{reality}.
The ideality can be thought, and there can be thinking about the
reality; but the reality itself cannot be thought; to think the
reality is to sublate it and transform it into ideality.

The immanent logic has shown excellently what ideality is, but it has
entirely forgotten what reality is.

When the finite, the existing positive, is to hold itself abstractly,
as the Real, over against the infinite as the negative, the Ideal,
then it is easy to see that the ideality must swallow up the reality.
For by having the infinite as its limit, the finite itself becomes
infinite; but thereby the infinite, which from its side receives the
finite as its limit, becomes precisely finite. In that the opposed
qualities have thus mutually negated one another, they have, by the
% --- fol. XIII ---
negation of the negation, at the same time affirmed one another, and
have thereby both obtained the same ideality: the finite's ideality
is its sublation and confirmation in the infinite; just so the
infinite's ideality is its sublation and confirmation in the finite.
But the reality has vanished.

The immanent logic has an altogether false conception of reality. In
its zeal to comprehend the object, it neglects to form an actual
concept of the actual object. The actual object is not an ideality
but a reality, a phenomenon of power, which no thinking is able to
produce. Not even the divine thought would be capable of thinking an
actual stone; for when the stone is thought, it is transformed into
ideality and ceases to exist.

It is with respect to this that the Logic of the Fundamental Ideas
has enforced the dualistic opposition of Knowledge and Power. If
Power is to pass straightway into Knowledge, then it is over with the
reality of existence; the things become semblance, and the actual
world dissolves into a metaphysical system of divine thoughts. If
Knowledge is to vanish into Power, then ideality is lost. Realism
sinks down into a crude materialism, which makes the Ideal a
semblance of the Real, the spiritual an empty after-glow of the
corporeal.

But when all the ideal content of being is elucidated in \emph{the
Idea of Knowledge}, and the real content of actuality passes into
\emph{the Idea of Power}, then it follows of itself that the unity of
the opposition must have its complete expression in the Idea of
Ideas, the all-sided Idea, which saturates thinking with ideality and
upholds the reality of actuality --- i.e.\ in \emph{the Idea of
Truth}.

\mainmatter

\chapter*{A. Knowledge.}
\markboth{A. Knowledge.}{A. Knowledge.}
% --- p. 1 ---
\section*{a). The Concept of Knowledge.}

To know something is, in indeterminate generality, to become
conscious of something. Here, then, are
presupposed\footnote{As regards the determination of philosophy's
first general presupposition, one must indeed distinguish between the
propaedeutic preparation and the systematic beginning. In the
philosophical propaedeutic the relation between the knower and the
object of knowledge is elucidated from the ever more concrete
standpoints of the doctrine of cognition, the doctrine of science,
and the doctrine of Ideas. (Cf. Philosophisk Propædeutik i
Grundtræk. Kbhvn. 1857). In logic, by contrast, the point of
departure is taken in pure, self-presupposing knowledge. (Cf.
Grundideernes Logik. I. Kbhvn. 1864. Side 1---11).}:
\emph{a knowing self, an object of knowledge, and consciousness's
unity of both}. In the apprehension of this unity what chiefly
matters is certainty, assurance, and insight. Knowledge is the
intellectual light in which selfhood-determinations and
determinations of the other are united and clarified in full
transparency.

\subsection*{α) Immediate Certainty.}

Through sensing and thinking the manifold content of existence
passes over into consciousness. But since thinking in its beginning
% --- p. 2 ---
is just as immediate as sensing, self-certainty and certainty of the
other are originally, though in different ways, immediate.

Immediate certainty is groundless, hence unfortified, and passes
involuntarily over into uncertainty. The objects alter: doubt arises
about the essence of the objects. The senses can deceive; doubt
arises about the testimony of the senses. Thinking develops
concepts; but the partially formed concepts stand in contradiction
with one another: doubt arises about the validity of thinking
(scepticism\footnote{On the sceptical dialectic cf. Phil. Propæd.
Side 152---154; Grundid. Log. II. 212---222.}).

In order that certainty may be won back through uncertainty, there
is required

\subsection*{β) Assurance.}

On the basis of immediate certainty, self-certainty was presupposed
without further ado, and there developed, in the clear light of
immediacy, a manifold knowledge of the other (the world of objects
and representations). But doubt came and darkened the clarity.
Doubt, which dissolves knowledge and confuses consciousness, begins
by undermining the certainty of the world of objects and ends by
attacking self-certainty. But in its attack on self-certainty doubt
dissolves itself and perishes in renewed certainty. Even if the
doubter could doubt everything (\textit{de omnibus dubitandum est}),
he could nevertheless not possibly doubt his own doubting. Doubt is
thinking, and the self is thinking (\textit{cogito ergo sum}); it
is, in its self-thinking, assured of its own being. But thereby it
becomes at the same time assured of the being of things.

From self-assurance follows the assurance of the subjective validity
of the original selfhood-forms (the pure intuitions, categories, and
Ideas) (Kant); but from this follows also the assurance of their
objective validity. If
% --- p. 3 ---
there could be otherness-forms, determinations of being in the thing
in itself (das Ding an sich), which by their nature stood in
conflict with the original selfhood-forms, then it would be not
merely impossible for us, but impossible in itself, to have a
knowledge of the essence of things; for not a single determination
can be taken up into knowledge without the knowing self having to
produce it. This holds not only of thoughts but also of
sense-impressions, not only of concepts but also of images and
representations.

Judging by immediate experience, one would indeed suppose that from
the world of things there were conveyed to consciousness countless
object-forms and property-determinations which did not in the least
concern its own essence. Yet upon some reflection it will readily
appear that nothing can be known except that which is determined in
knowledge. How should there be a knowledge of colour if the knowing
being did not itself have something to do with colour.

To appeal to visual organs and light-impressions is in this
connection of no significance, for the fundamental question is how
the organically conditioned light-impression is transformed into
consciousness of light, or how a colour-determination can become a
knowledge-determination. Could the light from without shine through
the organ into the soul, as through a prism it can fall upon the
wall of the chamber, then the colours would indeed be able to colour
consciousness, so that it became both yellow and blue.

\subsection*{γ) Insight.}

Through assurance the lost clarity is regained in the form of
clearing-up, grounding, and explanation, and corresponds to insight.
The knower's insight is double: it is an insight into the matter
with its members and relations, and it is an insight into the proper
essence of the thinking consciousness. Just as thinking of the other
is essentially self-thinking, so all thorough understanding of the
other is conditioned by self-understanding. Logical
self-understanding rests on the insight that the
% --- p. 4 ---
forms of reflection, judgment, and inference are universally valid
otherness-forms because they are original and necessary
selfhood-forms.

The knowing subjectivity is in its self-knowledge a reflecting
being. Immediately, reflecting is a mirroring, and presupposes a
mirroring activity and a mirror-image: pure knowledge is, in the
intellectual sense, pure mirroring. That the subjective activity is
mirroring makes it explicable that what is mirrored can remain
standing in the form of otherness without melting together with the
determining selfhood-form --- that colour-images, for example, can
be mirrored by consciousness and awaken consciousness of colours
without colouring consciousness. But so long as mirroring is
apprehended only as immediate reflecting, this reflecting will fit
otherness just as well as selfhood. In optics there are many mirrors
and many kinds of mirrorings; but logically understood, the original
reflecting is an ideal self-activity, a thinking. Things cannot
reflect, for they cannot think; on the other hand, not only can
their essence be thought upon, but this essence can be reflected in
itself when it is reflected in subjectivity: all reflection of the
other rests on self-reflection.

The knowing subjectivity is in its self-knowledge a judgment; herein
lies the ground of the judgment's form being an essential
otherness-form. Subjectivity is a judgment, that is, it is a subject
whose predicates are determined by the functions corresponding to
its judgments. ``The self is seeing, is hearing, is thinking''; by
the predicate ``seeing'' are determined all judgments about visible
things, by the predicate ``hearing'' about audible ones, by the
predicate ``thinking'' about thinkable ones. From this it follows
that in every possible judgment, whether the self for the rest
judges about the other or about itself, there must be a double
subject, the judgment's and the judger's, and likewise a double
predicate and a double copula: the otherness-form in the judgment is
infinitely reflected in the selfhood-form.

The knowing subjectivity is in its self-knowledge an
% --- p. 5 ---
inference; herein lies the ground of the inference-form being an
essential otherness-form. The inference arises --- so far as one
looks to what is typical --- in that the judgment's subject and
predicate, according to their difference of totality, form extremes
whose unity cannot immediately be indicated by the copula but must
be mediated by a middle term.

The selfhood-judgments declare, namely, that heterogeneous
predicates, belonging to different spheres, are posited in one and
the same subject. In seeing, the subject is the seer; in hearing, it
is the hearer; the seer is thus hearing, and the hearer seeing; the
sensing one is thinking, and the thinking one sensing. Raised above
the difference of the knowledge-functions, the self holds itself
fast as the singular. And yet the self is in itself a particular,
for sensing is something other than thinking; what is thought is not
sensed: the sensing one is thus not the thinking one, or: the
sensing and the thinking self are particular. Further, sight is
something other than hearing; what is seen cannot be heard, and what
is heard cannot be seen: the seer is thus not the hearer, or: the
seeing and the hearing self are particular.

Now by such a division the self would be split asunder if the
division were not a self-division that let itself be equalized in
the universal. But the self is itself this equalization.

The self divides itself into its functions; but these functions,
however heterogeneous they may be, are yet all selfhood-functions,
and, as selfhood-functions, functions of one another mutually.

If in sensing there were not a hidden thinking, the sensing self
could distinguish nothing, and consequently could not sense: sensing
is thus a function of thinking. Conversely, if in thinking there
were not an implied sensing, there would be no element of
representation, and the thinking self could distinguish nothing, and
consequently could not think: thinking is thus a function of
sensing. One distinguishes in
% --- p. 6 ---
sensing by the help of thinking, and one distinguishes in thinking
by the help of sensing. That there is distinguishing is the
universal; that which under the universal distinguishing is
separated and distinguished is the particular. The distinguisher is
in its distinguishing equal to itself: the self that distinguishes
itself in everything and everything in itself is thus the particular
because it is the universal.

Apprehended immediately for itself, the self is the singular. In
this its singularity it is a fixed point, an atom, from which
activities go forth in all directions. To this attaches the
representation of an existent thing, a soul-real (Herbart); but such
a representation makes the self selfless.

Apprehended immediately in the form of particularity, the self is a
swarm of selfhood-points: it is the centre of the whole circle, but
every point in the circle is in turn a centre. This monstrous
contradiction stems not from particularity, but from the immediacy
in which the particular is held fast. The dissolution of immediacy
in infinite reflection is the solution of the contradiction. The
self is reflected in its singularity, for its activity as a singular
is, by means of the particular, a self-activity. It is reflected in
its particularity, for its particular activity is, by means of its
singularity, a self-activity: self-activity $=$ self-activity is the
universal in which the self consists.

Knowledge is not activity, but the light in the universal activity,
the infinite reflecting and mirroring. The universal activity, an
activity of activity into the infinite, is in its formal generality
an empty concept. The essence that is universal in the absolute
sense is thus self-essence only in that the singular and the
particular, each for itself in mutual opposition, is self-essence.

Where then is the self in this essence? It is in its own inference.
When the singular and the particular are posited as extremes, the
middle is the universal (S---U---P); when the singular and the
universal are posited as extremes, the middle is
% --- p. 7 ---
the particular (S---P---U); when finally the particular and the
universal are posited as extremes, the middle is the singular
(U---S---P).
% The Danish sigla are E---A---S, E---S---A, A---E---S
% (Enkelte/Almene/Særskilte); rendered with the English initials
% S(ingular)/U(niversal)/P(articular) in the same order.
The self has subsistence in a system of inferences: the
inference-unity, the absolute central unity, is itself the eternal
beginning-unity.

From this standpoint it will now be possible to attain insight into
the relation between the logical fundamental forms of selfhood and
of otherness.

In the essence of knowledge they are the same forms; the opposition
rests solely on reflection's position toward the immediate. In the
selfhood-forms immediacy is completely dissolved; but in the forms
of otherness the immediate constantly returns and interrupts
reflection. Otherness cannot reflect itself; it must be reflected by
selfhood.

As object for the self, that which is is itself a something for
itself. That it is something for itself points to selfhood; and yet
something is for itself only what it is in its relation to an other.
But something passes over into an other only so far as it is thought
together with the other; its transition is no self-movement --- the
movement is in the thinking subjectivity. Something reflects itself
in an other only in that the self reflects upon something and other:
in becoming and alteration the otherness-reflection shows itself
expressly as the intellectual product and counter-image of
self-reflection.

Otherness-being is becoming. If that which is is reflected in its
own becoming, and conversely: then immediacy is sublated; then that
which is is essence, and that which becomes, semblance. As essence
and semblance reflect themselves in one another, it looks as though
the essence itself thought its semblance, and the semblance its
essence; but the thought-movement goes on in the thinker. The
reflection of semblance and essence is self-reflection's own
reflection of the other. That the reflection-form, essential in
itself, is an otherness-form is seen also from this, that the
essence's reflection into itself is posited as thing, and its
reflection into an other
% --- p. 8 ---
as property: the thing itself is selfless, and the property a
quality of the selfless.

But the selfless must necessarily acquire a semblance of selfhood;
otherwise the otherness-forms could not possibly become
knowledge-forms. The knowing self then knows that the concept of the
thing's being is a judgment. The thing now becomes subject, and the
property its predicate. The subject is its predicate: it looks in
otherness quite as in selfhood. But in the essence of otherness the
selfhood is semblance; the thing, the selfless subject, has no true
subjectivity; the subjectivity lies on the surface, for it lies in
the properties.

By making the world of things into a world of judgments, the judging
self knows how to hold things together and to comprehend the unity
of connection in a system of inferences. Unity of connection is
determined by relation, relation by reflection, reflection by
judgment, judgment by inference: the articulated reciprocal
determinations of the forms of selfhood and otherness yield the
fundamental moments in the concept of knowledge.

\section*{b) Finite Knowledge.}

According to the concept of knowledge, self-knowledge and knowledge
of the other must necessarily presuppose each other into the
infinite; but in the special development they part company, and
herein finitude is seen. The distinguishing mark is that the
knowledge-content which saturates the selfhood-forms and is
necessary for their logical unfolding is poor, empty, and formal in
comparison with the content that fills the forms of otherness with
their outwardly infinite wealth. The strict, systematic deduction of
the original I-hood propositions, as it has been carried through in
subjective idealism (Fichte), is a telling example of the a priori
poverty of pure self-knowledge. In contrast, the doctrine of
experience has such a superabundance of content-determinations and
objective knowledge-forms that
% --- p. 9 ---
along this path there can be no question of demonstrating the
necessity with which the facts of the objective world proceed,
member by member, from the subjective fundamental demands of the
I-essence (nature-philosophical mysticism; Schelling).

In order to make the relation of finitude clear, we note especially:
first the one-sidedness of a priori knowledge, then that of
empirical knowledge, and thereby come to see the necessity of the
logical dissolution of finite knowledge.

\subsection*{α) A priori Knowledge.}

Formal logic's conception of the principle of identity and
contradiction finds its most adequate application in the science
that is precisely suited to make intelligible the finitude of a
priori knowledge, namely mathematics. The equation: A $=$ A, is in
the objective sense what the self-proposition: I $=$ I, is in the
subjective; likewise the proposition: A is not Not-A, in the form of
otherness, has as its presupposition the proposition: I am not
Not-I, in the form of self-limitation. Here is immanent reflection
of the subjective and the objective.

Further, a hypothetical proposition such as this --- if
$\frac{a}{b} = \frac{c}{x}$, then $x = \frac{bc}{a}$ --- can serve
as an example that objective relation is grounded in subjective
reflection. From the antecedent the consequent must necessarily
follow; the relation between a, b, c, and x is so far objective. But
the magnitudes have not themselves set themselves in relation to one
another; a has not itself divided itself by b, nor c by x. The
relation is posited by the reflecting subjectivity. In the quotients
$\frac{a}{b}$ and $\frac{c}{x}$, a is reflected in b, and b in a, c
in x, and x in c; the numerator presupposes a looking to the
denominator, and the denominator to the numerator: that is
subjective reflection. If an equation such as $\frac{dy}{dx} = c$ is
chosen, the numerator is no immediate term, but a reflex of a
magnitude vanished in vanish-
% --- p. 10 ---
ing, and the same holds of the denominator dx; as the two reflexes
are reflected in each other and form the quotient, the reflection
becomes objective; but it is only objective because it is
subjective.

Could the mathematical development now be continued in such wise
that to every new objective relation there corresponded a new and
fruitful insight into subjective reflection, to every new theorem a
new and fruitful cognition of subjectivity's original self-positing:
then we should here have a method (\textit{methodus mathematica})
which was in an eminent degree suited to sublate the finitude of a
priori knowledge.

But this reckoning does not suffice. It is true, indeed, that
propositions such as, A is A, A is not Not-A, belong so intimately
to general self-thinking that their validity is evident to every
thinker. But whence comes it that equations such as
$\log A^{x} = x \log A$,
d.\,$\sqrt{a^{2}-x^{2}} = \frac{-\,x\,dx}{\sqrt{a^{2}-x^{2}}}$,
even when the notations are explained one by one, cannot be said to
belong to the original comprehension of the self, are not
indispensable for logical self-understanding? Nay, whence comes it
that mathematics cannot be deduced from metaphysics, or conversely?
It comes from this, that a priori knowledge, on account of
finitude, parts into two one-sided directions: a logical-dialectical
one, which immerses itself in self-reflections, and an
exact-operating one, which loses itself in reflections of the
other\footnote{On the difference between reflection-concepts and
operation-concepts cf. Mathematik og Dialektik. Kjøbhv. 1859.
S. 40---56.}.

\subsection*{β) Empirical Knowledge.}

The a priori otherness-forms that are articulated in the pure
analyses of magnitude rest solely on phantasy-intuitions and
representations, and the represented otherness falls, so far, still
within selfhood.

But when it is a question of a knowledge of existing things,
% --- p. 11 ---
the contradiction enters that the objects within knowledge
themselves fall outside knowledge. The finite consciousness has,
namely, a world of objects \emph{outside} itself; but it cannot
possibly have a knowledge of these objects except by taking their
being up into itself, and thus having them \emph{within} itself.
This contradiction, hidden in the foundation of empirical knowledge,
conditions the manifoldness of cognition, but at the same time
betrays its finitude under the form of contingency.

The contingency is subjective, for the universally necessary
otherness-forms in which the objects must be apprehended can indeed
be had in consciousness without knowledge that the special objects
really exist. That an existent object must always correspond to some
concept or other, that is, to some filling-out or other of a priori
forms, does not sublate the contingency, since such a filling-out
can take place in countless ways: that many and different things
bear the same name points precisely to the differences being
regarded as contingent. The empirical concept to which the name
answers repeats itself in a multitude of corresponding specimens; on
the basis of the concept the specimen is described, but the
pictorial elements of the description indicate the contingency.

In relation to the individuals and the immediately individual, the
existent concept is the species, and the nearest universal concept
of the particular species is the genus. That an individual with one
or more marks of a certain species must have all the marks belonging
to the species is inferred by \emph{analogy}; that essential
determinations which are decisive for some individuals belonging to
a certain species or genus must hold for them all is inferred by
\emph{induction}; but since not only the individuals but also the
concepts are empirical, both induction and analogy are burdened with
a moment of contingency that can transform certainty into
uncertainty.

From a general standpoint empirical knowledge can indeed attain
reliable certainty; but the particular cases
% --- p. 12 ---
constantly cross the universal and call forth uncertainty. The
ground of assurance is negative: ``here is no case that breaks with
the universal established by the observations; ergo this universal,
determining the cases, is a law.'' But the universal thus
established is nevertheless, in its positive form, empirical. To the
methodically continued observations the world of experience
constantly presents new sides, new relations, new properties, new
laws.

The contingency is objective, for it is not merely we who must
alter, correct, and improve our concepts of things; the concepts
themselves (species and genera; Darwin) alter by degrees, in that,
as existence-concepts, they alter the existential expression with
the alterations of the individuals. The contingency is objective,
for the existence of the objects is independent of our subjective
knowledge of them. We could not sense the objects if the objects did
not sensibly exist; but the sensible objects can very well exist
without our sensing them. We form for ourselves concepts that are
abstracted from the objects, but the objects themselves are
indifferent to our abstracting and comprehending. It lies in the
independence of the objects to have our knowledge of them outside
themselves as a contingency.

For our finite knowledge, then, the world of things stands as a
barrier, an impenetrable other. We cognize, says the doctrine of
experience, the properties of things, but not the things themselves;
we cognize the laws, but not the essence. Mechanics does not know
what force is; chemical physics does not know what electricity is;
physiology does not know what life is. True empirical knowledge has
everywhere taken up into its foundations the determination:
not-knowledge.

\subsection*{γ) The Dissolution of Finite Knowledge.}

The contradictions with which finite knowledge is burdened on all
sides gather at last into the one: omniscience is necessary and yet
incomprehensible.
% --- p. 13 ---
That omniscience is necessary shines forth from the original essence
of knowledge. For since all determinations of existence are, in the
logical sense, knowledge-determinations, and all
knowledge-determinations functions of the knowing self: nothing
whatever can be known except that alone which is produced by
knowledge.

The contradiction is now this, that finite knowledge everywhere in
the actual world finds before it object-forms which neither it
itself nor the object has produced. All outward things exist in
space. But what is space? An extension produced by intuiting,
existing for intuition. In intuiting the things in space, we intuit
them in a previously intuited form (objective space); and the same
holds of time.

By reflecting on the objects existing in space and time, we come to
a similar result. For us, with our finite knowledge, the intuitable
objects exist before we intuit them. But intuitable is only that
which is produced by intuiting. Our intuiting is determined by the
intuitable, that which is already produced. How has the object
become intuitable? Solely in that its existence is posited in an
intuitable form, i.e.\ in a form possible for intuition; but an
object-form possible for intuition is actual only in virtue of an
original intuiting that forms the object.

For the original intuiting, the object itself is reflected as image,
and the image as object. In the form of the image, however, the
object-reflection is superficial; the essential content has its
exhaustive form in the concept. For our knowledge it is the object
that determines the concept; but in the matter itself it is the
concept that determines the object. In order to determine the
object, the concept must beforehand be determined by a comprehending
independent of our concept, hence by a knowledge independent of our
knowledge.

The intuiting, thinking, and comprehending of omniscience, in which
existence should have its original ground, is indeed
% --- p. 14 (sentence and subsection conclude here) ---
incomprehensible from the standpoint of finitude (simultaneous and
successive knowledge; intuitive and discursive
knowledge;\footnote{Cf. Grundid. Log. I. S. 230---63.} knowledge's
relative and absolute limit); but as omniscience in represented form
is incomprehensible to finite knowledge, finite knowledge becomes at
the same time incomprehensible to itself: in order to comprehend
itself, finite knowledge must recognize its dissolution in the
infinite.

\section*{c) The Infinite Knowledge.}

So far as the original knowledge is immediately apprehended as
omniscience, it looks as though it were only a boundless extension
of man's restricted knowledge: ``Man knows something, the learned
man knows much, but God knows everything.'' This conception is
unscientific. The concept of infinite knowledge cannot be formed by
sublating the limit of the finite; for when the limit falls away,
everything runs together into one. The fundamental difference
between finite and infinite knowledge rests on a thoroughgoing
difference between \emph{restricting} and \emph{indwelling} limits.

Finite knowledge has that which limits it outside itself, because
the finite self has its other, its actual opposite, outside itself;
our self-knowledge is in its finitude separated from our knowledge
of the other; in its formal circular movement it has only the shadow
of the infinite, not infinity itself. The divine self, on the other
hand, has its other, its actual opposite, in its own essence: for
this self-knowledge the immediate \emph{outside} is essentially an
\emph{inside}. The infinity of knowledge shows itself in reflection,
in judgment, in inference.

\subsection*{α) The Infinity of Reflection.}

The sensible is the immediate other for sensing; likewise the
thinkable is an other for thinking. If now this sensible and
thinkable other is apprehended as outward limit, as
% --- p. 15 ---
barrier, the sensing-thinking consciousness gets a world of
sensible-thinkable objects outside itself.

Dazzled by this outside of immediacy, finite knowledge then assumes
that the sensible stands beforehand fully finished and merely waits
to be sensed. This waiting, then, signifies no lack or essential
want. If the sensing comes, well and good; the sensible takes no
harm from it, suffers nothing under it: it willingly permits the
senser to apply his apprehension, for it is indeed there to be
sensed, it is indeed sensible. But if the sensing does not come,
that too is well; the sensible has all the same its complete
determinations of being, its pronounced actuality-forms in itself.
And as it goes with the sensible in things, so it goes also with the
thinkable. This way of seeing is a telling proof of the restriction
of finite knowledge.

Sensible and sensing, thinkable and thinking, are correlate
concepts. What, indeed, is the sensible for itself? It is an
indeterminate that must let itself be determined --- by sensing. If
the sensible is held fast in all its indeterminacy, it is not an
existent; for everything existent is a determinate being: what is
still indeterminate does not exist in actuality, it is only a
possibility. The sensible is a possibility, namely for sensing, a
possibility of sensing, just as the thinkable is a possibility of
thinking. Outside all sensing and thinking, the objects are, not
merely for us but also in themselves, indeterminable magnitudes,
pure possibilities without any actuality.

But now the converse holds as well, that sensing for itself is only
a possibility for the sensible, and thinking for itself a
possibility for the thinkable. Where there is nothing sensible there
can be no sensing, and where there is nothing thinkable there can be
no thinking. To sense a nothing is impossible; and when thinking
thinks its nothing, its
% --- p. 16 ---
negation, this nothing is itself, by virtue of the word, a thinkable
that acquires existence for thought. The two possibilities, that of
sensing and that of the sensible, corresponding to that of thinking
and that of the thinkable, are in knowledge one and the same
possibility; they are actualized through each other; they stand and
fall with each other: they must either both become actual or both
impossible.

If now it is incontestably certain that all sensing must have a
sensible, and all thinking a thinkable, for its possibility; if it
is evident from the matter itself that the sensible can become
actually determined only by the determining sensing, and the
thinkable only by the determining thinking: then this is to say, in
other words, that the sensible-thinkable objects are indwelling
limits, real-negations in the proper essence of sensing and
thinking. As real-negation the sensible-thinkable steps forth with
object-like immediacy, and so far there is here an actual
\emph{outside}; but since this immediacy, under all turns, in all
alterations of form, is only the real-negation's own appearing, the
actual outside is everywhere, in knowledge's infinite reflection,
essentially an \emph{inside}.

\subsection*{β) The Infinity of Judgment.}

Sensing is not thinking, and the sensible is not a thinkable; and
yet sensing is possible only through thinking, the sensible possible
only by means of the thinkable: sensing is at bottom thinking, and
the sensible thinkable. Both determinations --- both that sensing is
not and yet is thinking, and that the sensible is not and yet is
thinkable --- are taken up into the infinite reflection; but in
order to solve the contradiction contained therein, reflection must
be expressly determined by the judgment.

The categorical judgment: ``The subject is (essentially) its
predicate,'' receives its fullest expression in the concrete form:
``The individual is its species.'' All determinations belonging to
the individual are sensible; all marks answering to the species as
% --- p. 17 ---
species, on the other hand, thinkable: the individual is sensed, the
species is thought; thus the judgment expresses at once the
difference between sensible and thinkable and at the same time the
unity.

The finite knowledge, restricted by its \emph{outside}, cannot
however come to clarity with itself concerning the relation between
subject and predicate. It looks as though the individual, the
sensible, were properly the first and original, and the species, the
thinkable, on the contrary the derived; for the judging subjectivity
must indeed set out from a sensible apprehension of the individual
and then abstract the species from the individual. Thus the existent
individual is \emph{outside}, the abstracted species-concept on the
contrary \emph{inside}, consciousness.

But if this relation is the properly true one, then it must be the
individual that determines the species, while thinking on its side
maintains that it is the species that determines the individual.
Sensing asserts that the individual is the original, the species the
derived; thinking asserts that the species is the original, the
individual the derived: here is tension and strife between sensing
and thinking.

In the infinity of the judgment the contradiction is solved in that
the two factors, the sensible and the thinkable, are both taken up
as moments in one and the same essential fundamental factor, namely
in the selfhood-act whereby the judging subjectivity has originally
posited the judgment's content in the judgment's form.

It is, then, just as little the abstract concept, the species, that
produces the individual, as it is the immediately existent
individual that yields the species-concept. The individual's
existence has the species as its ground, so that the species, the
content of thinking, grounds and determines the object of sensing,
the individual. And why? Because the sensible is from the ground up
contained in the thinkable, because thinking in the original
judgment-act rests precisely on sensing, just as conversely sensing
on thinking.
% --- p. 18 ---

In the finite judgment-act the moments fall apart from one another:
the individual, the judgment's subject, is found for itself as a
sensible object; then thinking steps in, abstracts the
species-marks, and determines the predicate; then it looks as though
the individual fell outside the judgment and remained indifferent to
the whole judgment-act; then the judgment-act has significance only
for the judging subject, not for the subject in the judgment. It is
a parcelling-out that characterizes finitude.

In the infinite judgment-act, on the other hand, the moments are
gathered in one comprehending into one concept. In this
knowledge-concept sensing is infinitely articulated by thinking, the
sensible by the thinkable, the outward by the inward: the piecework
is sublated, but the particularity of the members, the difference of
the moments, are all preserved in the unity of transparency.

The disjunctive judgment: ``A is either B or C,'' is in formal logic
a schema, an empty form, which can take up from without many kinds
of content. It receives its fullest content through the
determination of the relation between species and genus. The genus A
is either the species B or the species C, and so on. Formal logic,
which places the judgment-act down in the finite subjectivity, now
conceives the relation between the genus and its species in such
wise that the genus-concept has the lesser content because it has
the greater extension, while the converse holds of the
species-concept. According to the immanent logic it is, on the
contrary, the genus which contains the species within itself. The
genus-concept is then determined as the content-rich concept, the
concrete universal concept, in which the species with their
one-sidednesses complete one another. This conception has certainly
aimed at the form of infinity; but the unovercome finitude shows
through, for the form cannot logically master its content.

As the concrete-universal, namely, the genus-concept should
% --- p. 19 ---
itself comprehend its species-concepts; it should, without giving up
its own content-rich unity, part with the species and set them so
exclusively against one another that each species for itself
excluded the rest. But that is impossible. If the separation of the
species is in earnest, the genus-concept loses its concrete content,
and formal logic proves right. If the genus is absolutely to have
the same content as the species, this can happen only in one of two
ways: either the particularity of the species must be equalized in
the generality of the genus, and then the species perish in the
genus; or else the universal as genus-concept must exhaust its
content in the particular species, and then the genus perishes in
the species. The judgment demands a judgment-act whereby both things
happen, but in such wise that the contradiction is solved.

Such a judgment-act no genus-concept can consummate. The
genus-concept is subject in the judgment; but it is not judging
subject. The infinity of the judgment is grounded in the judging
subjectivity's infinite judgment-act: 1) the content laid out in the
predicate is and remains the subject's content; as the judgment's
subject, the genus-concept reflected into itself has the whole
undivided, undiminished content of the species-differences in its
own universal form; for precisely in this universal form it is
subject. 2) The ground-division in the predicate is carried through
in such wise that each member has in particular form the whole
subject-content; each species is a particular expression of the
whole genus; 3) finally, the opposition of subject and predicate is
completely preserved in the unity determined by the judgment-act.
Thus the judgment's members and moments are posited apart from one
another, and yet they are in one another: every outside is an
inside.

\subsection*{γ) The Infinity of Inference: the Idea of Knowledge.}

As the judgments are, so must the inferences drawn from them also
be. Just as the difference between finite and infinite knowledge
has, as regards the judgment, been determined by the
% --- p. 20 ---
judgment-act, so it is now, as regards the inference, to be
determined by the inference-act.

The immanent logic has, by dialectical presentation of the concept's
trilogical determinateness as: \emph{the Universal}, \emph{the
Particular}, and \emph{the Singular}, certainly given the inference
a form that suits the absolute content, hence the Idea. But the
ambiguity of the form-development lies hidden under an ambiguous
conception of the inference-act. For an inference can no more infer
itself than a judgment can judge itself, or a concept comprehend
itself. The self that comprehends itself in all concepts is in its
self-comprehending raised above every concept, in its self-judging
above every judgment, in its self-inferring above every inference.
But precisely for that reason this original, unoriginated self is
the essential source of all concepts, all judgments, all inferences.

So long as one, as in formal logic, lets the psychologically
restricted subjectivity carry out the inference-act in all
inferences, or, as in the immanent logic, abstracts from finitude by
abstracting from the inferring subjectivity, so long can one go on
reducing and transforming the inference-forms as much as one will:
one never comes to clear insight into the Idea of Knowledge.

The old strife about genera and species (γένη and εἴδη), whether
they are Ideas or abstract concepts, cannot be settled in finite
knowledge. But the finitude of knowledge persists so long as it has
not yet been brought to full consciousness that the infinity of the
inference is the Idea's logical form, and that this infinity rests
decisively on the infinite inference-act.

The individual in sensible immediacy is indeed no Idea; but a
concept in abstract intellectual form, whether for the rest it be a
species-concept or a genus-concept, is no Idea either. That pure
sensing cannot apprehend the Idea is certain; but that pure thinking
--- das Denken
% --- p. 21 ---
des Denkens im allgemeinen Denken --- likewise cannot apprehend the
Idea is also certain. To regard the substantial universal forms as
self-valid realities and the individuals as fleeting phenomena
(\textit{universalia ante rem}) conflicts with the Idea of Knowledge
just as much as to place the whole reality in the individual and
reduce the universal forms to empty names (\textit{nomina, flatus
vocis}). The original relation between the universal forms and the
individuals can be determined solely by the inference-act.

If it depended merely on setting up an inference-schema such as
this: the individual is the genus by means of the species
(S---P---U), the species is the genus by means of the individual
(P---S---U), the individual is the species by means of the genus
(S---U---P): then Idea-cognition would be not only the easiest but
also the most superficial of all scientific cognitions.
% Danish sigla: E---S---A, S---E---A, E---A---S, with E = Individ
% (det Enkelte), S = Art (det Særskilte), A = Slægt (det Almene);
% rendered with the English initials S/P/U as at p. 6 f.
The task first comes into view when one reflects on what such an
inference-schema essentially means.

The schema set up means, namely, that of the inference's three
members --- individual, species, and genus --- none is the first and
none the last, but that all three must have the same originality.

The sum of sense-determinations gathered in the unity of the
individual is determined beforehand by the totality of
thought-determinations contained in the species; but this
thought-content is in turn a modified form of the universal content
of the genus-concept. So the formation of the genus-concept would
accordingly be the unconditioned first. But the formation of the
genus-concept presupposes precisely the species-concept, and this in
turn the individual; so the individual would accordingly be the
unconditioned first after all. Here is the contradiction. The
individual cannot be formed before the species and the genus are
formed; for, \emph{seen in the Idea}, \emph{the individual is the
immediate inference-unity of species and genus}. The species cannot
be formed before the individual and the genus are formed; for,
\emph{seen in the Idea}, \emph{the species is the
% --- p. 22 ---
reflected inference-unity of genus and individual}. Finally, the
genus is \emph{itself, determined as Idea, the comprehended
inference-unity of species and individual}, and must expressly have
the two extremes as its premises.

It is then surely evident that the inference-system which is to
express the Idea's content in the form of knowledge must spring from
one single threefold inference-act, and that no psychologically
limited subjectivity is in a position to carry out such an act.

Logical insight into the infinity of the inference, and a
corresponding formal insight into the Idea's original essentiality
--- hence a knowledge of the Absolute --- is made possible by the
dialectical dissolution of finite knowledge in the infinite; but
such a knowledge of the Absolute is infinitely far from being an
absolute knowledge.

To have a knowledge of the Absolute that is true in the formal
respect is to have an absolute assurance that all determinations of
being are originally knowledge-determinations, that the light of
knowledge is the essence of things, that the content of the whole
existing world is grounded and borne by infinite, i.e.\ by divine,
knowledge.
% p. 22 ends the text of A. Knowledge; a short centered rule follows
% below the last line in the print (carried as a comment here, as in
% the transcription).

\chapter*{B. Power.}
\markboth{B. Power.}{B. Power.}
% --- p. 23 ---
\section*{a) The Essence of Power.}
% (no stop after »a)« in the print, unlike »a).« in ch. A)

Although all being comes in under knowledge, there is nevertheless
in actual existence an other, namely the immediate object-being,
which forms an opposition to knowledge; this other, this
opposition-determination, has from the standpoint of infinity
hitherto been apprehended only as indwelling limit-determination, as
real-negation.

If the otherness-limit, the real negation, which makes the
knowledge-determination into a reflected object-determination, did
not --- precisely on account of otherness's own essence --- part
itself into a manifold of negations and object-reflexes, the
knowledge-concepts would melt together into an empty semblance, and
the unbroken light in the Idea of Knowledge would be differenceless,
like pure darkness.

But a negation is impossible where there is nothing that can be
negated; every negation is originally a negation of something
positive: that the negation of the negation yields a positive
determination means that the positive, whose immediacy was sublated
in the first negation, returns with the second. The indwelling
real-negations of the Idea of Knowledge thus presuppose a
corresponding manifold of immediate object-positions --- an
expression of the essence of power.

% --- p. 24 ---
\subsection*{α) Negative and Positive in the Essence of Power.}

Were it possible to hold fast the negative for itself without the
positive, and conversely, one would now be able to separate
knowledge and power thus: the essence of existence --- the
thought-content reflected in negative determinations, negatively
affirmed through negations and negations of negations --- comes in
under knowledge; actual existence --- the object-world existing in
positive determinations, unfolding itself in immediate positions ---
comes on the contrary under power; in other words: the world of
concepts (\textit{mundus intelligibilis}) has its subsistence in
knowledge; the material world of things (\textit{mundus
sensibilis}) has its subsistence in power. In this way it would then
come to look as though power were joined on to knowledge, and
actuality to essence.

That such a compounding of knowledge and power is hollow and
contradictory is understood of itself.

The positive is indeed qualitatively different from the negative;
but this difference is precisely grounded in the unity of the
essence. The same essence that negates must also be able to posit,
for otherwise the negation could not sublate the position, could not
be essential negation; conversely: the same essence that posits must
itself be able to negate, for in the contrary case the negation of
the negation could not possibly become any true and actual position.
This the immanent logic has developed to convincing clarity, and all
later attempts to hold the absolutely positive sharply separated
from the negative (Herbart) have shown themselves invalid and
fruitless.

But then the immanent logic has in turn committed a great blunder,
in that it has without further ado let the negating and positing
essence be absorbed both in its negations and in its positions. If
the negating essence is completely absorbed in its negation, then
the negative negates itself, independently of the position; if the
positing essence is immediately absorbed in its position, then the
positive posits itself,
% --- p. 25 ---
independently of the negation; in this way the one and the same
essence would, in spite of the principle of immanence, split into
two independent essences, a negating and a positing one. If the
negating essence is itself to be positing, it cannot, on account of
the position, be absorbed in the negation, and just as little, on
account of the negation, be absorbed in the position; if it is to be
sunk in the infinite alternation of positions and negations, then,
strictly taken, it does not even become the ground of all that
becomes; it becomes process of process, real becoming pure and
simple, a becoming that is not essence.

If the opposition of negation and position is to be thinkable as
proceeding from the same essence, then this all-negating,
all-positing essence must be a self-essence raised above its
negations and positions; only such a self-essence is capable, by
sublation of barriers in the other (the immediately positive), of
affirming itself through negations in knowledge, and, by negation of
negations, of affirming its other outwardly through positions in
power. The outward-going affirmation, the affirmation of the other,
is then actuality's measure of the self-affirmation, the
inward-going affirmation: the indwelling capability of this
double-sided affirmation deposits
% (as in the Danish, the sentence runs grammatically into the
% following run-in head as its object)
\subsection*{β) Extremes in the Essence of Power.}

Power is essence, be it noted, in the form of actuality, developed
to absolute necessity. Categories such as: force, possibility,
condition, and existence form the development's abstract logical
basis. Force, regarded as possibility, is condition for the
manifestation, and the manifestation is the existence that results
where the condition takes place; but necessity is the activity
which, by completing the reflection and sublating the condition,
leads the existent back to immediate being. The necessary is indeed
necessary by means of an other, which is its condition, and this
condition is itself existing or independent; but since the necessary
does not itself come
% --- p. 26 ---
to existence before the condition's existence is sublated, the
necessary is just as much unconditioned, or: it is because it is. In
this unconditionedness the necessary makes itself immediately felt
as power. --- From this standpoint the concept of power is
apprehended in the immanent logic.

But what has cognition now gained by such a conception? A
dialectical demonstration of the place which, according to
scientific usage, belongs to the category: power, among the other
logical categories; nothing else. But what was to be attained here
is something quite different; here, namely, an insight was to be
attained into how actual being relates itself to thinking, or, still
more definitely, how the existing thing relates itself to its
concept.

The problem of \emph{thing} and \emph{concept} the immanent logic is
so far from having solved that it has not even been able to set it
up. But that here there really is a problem can be seen from this,
that here there is a contradiction.

The contradiction is that thing and concept both are the same and
yet are not the same. They are the same; for there can no more be
thought any decisive peculiarity in the thing that should not belong
to its concept, than there can be assumed in the thing's concept any
determination that does not belong to the thing; and yet they are
not the same, for the thing is sensible, but the concept is
non-sensible. They are the same, for concept and thing have the same
content; and yet they are not the same, for the concept, the
universal, has an infinitely greater extension than the thing, the
singular (one concept can be in many things). They are the same, for
the existing thing is the concept's existence; and yet not the same,
for the existing thing is no concept, and the concept no existing
thing (a stone is no concept, and the stone's concept is no stone).
One can crush a stone, but not the stone's concept; one can boil a
plant, but not the plant's concept; one can ride on a horse, but not
on
% --- p. 27 ---
the horse's concept. When the thing is posited, its concept is
sublated; when the concept is posited, the thing is sublated, or:
position of the concept is negation of the thing, and position of
the thing is negation of the concept.

So there is required here an essence-doubling which from the ground
up makes the opposition of thing and concept intelligible, in that
it at the same time makes the unity of the opposition intelligible.

For the concept, the first extreme, we have knowledge; for the
thing's immediate reality, the second extreme, power; the middle,
the unity of the extremes, is self-unity: the knowing self is
powerful, the powerful self knowing; the powerful knowledge is, in
the unity of the inference, an ideal capability, the knowing power a
real one.

The inference-unity, the original capability from which the extremes
go forth and in which they subsist, is an active unity; its being is
activity, and its activity an infinite converting of concept into
existence, of existence into concept.

The conversion takes place by a leap, for the concept is
qualitatively different from the thing: the concept cannot be
condensed into thing, and the thing cannot be thinned out into
concept. But it is a leap whereby nothing is leapt over; for power
without knowledge, blind power, is just as much a semblance as
knowledge without power. The concept that has absorbed power in the
form of knowledge is a powerful concept, a concept that is capable
of determining the thing; conversely, the thing determined by the
concept is an expression of knowledge in the form of power, a
concept uttered in existence.

Since now the knowledge impregnated with power, by its negative
form, deprives power of all immediacy, transforms it into logical
capacity, into intellectual capability: it makes power invisible,
posits the concept and sublates the thing. The power impregnated
with knowledge-determinations, on the other hand, interrupts the
reflection; it sublates thinking by its immediate,
% --- p. 28 ---
positive form; it takes the concept up into real capability and
brings the capacity's logical tension to antilogical discharge ---
that is, it posits the thing and sublates the concept.

\subsection*{γ) The Logical-Antilogical in the Essence of Power.}

The concept can be thought, but the thing cannot be thought; the
concept can be uttered in words, but the thing's existential
determinateness is unutterable: the concept is logical, the thing is
antilogical.

It is indeed assumed without further ado that it is immediately the
thing that is thought, and that the word exactly expresses what the
thing is; but that is a deception which the logical judgment
occasions. What now is This? This is This, and yet the pointing
designation can expressly show that This is not This; for This, for
example, is a house, but This is a horse, and This is in turn
another This. By pointing to the object the word points away from
the word, the thought away from the thought: this is the
antilogical.

The antilogical stands over against the logical as the object stands
over against its concept; but just for that reason the antilogical
can no more be either illogical or in contradiction with the
logical, than the object can be either concept-less or in
contradiction with its concept. The unity is that the object and the
concept have, as far as the essence is concerned, the same
determinations; but the opposition is that the object-determinations
are antilogical, the concept-determinations on the other hand
logical.

That the same essential determination which in the form of thinking
is logical can in the form of actuality be antilogical, every
judgment about existing things proves. The existence-judgment can no
more dispense with the antilogical than it can dispense with the
logical. If the antilogical is placed in the subject and the logical
in the predicate, the thing is judged to be concept, the for-itself
unthinkable to be thinkable, and the judgment is logical (This is a
rose). If conversely the logical is placed in the subject and the
antilogical in the predicate, then
% --- p. 29 ---
the concept is judged to be thing, the thinkable to be unthinkable,
and the judgment is antilogical (The rose, the concept, is This, the
specimen).

The immanent logic treats the actuality-concepts without further ado
as pure knowledge-concepts; that is a great error. The
actuality-concepts are power-concepts, recognizable by this, that
the antilogical everywhere breaks forth out of the reflection and
exchanges thought for fact. Categories such as: actuality,
substance, cause, etc., are logical only so far as they are at the
same time antilogical. Actuality is a concept, but the concept is no
actuality; substance is a concept, but the concept is no substance;
cause is a concept, but the concept is not cause.

Power's antilogical discharge, the transition of possibility into
actuality, takes place with necessity; but this necessity is, by
means of the antilogical itself, a contingency. A line, for example,
is to be drawn. The line's concept is possibility of countless
lines; but the line that is drawn is expressly this determinate one.
Let one draw whatever line one will, it becomes each time a
\emph{This}. This is accordingly necessary; but since it could just
as well have been any other whatever, precisely \emph{This} is
contingent. So with every existing case.

If reflection is made on a series of cases, of which each preceding
one makes its successor necessary, we get the contingent necessity
more closely determined as \emph{conditioned} necessity. But this
logical concept still decides nothing concerning the actual cases.
The general series does not exist; each existing series is a
\emph{this} series, each member in this series is a \emph{this}
member; but this existing member, though grounded in the whole
connection, is an antilogical discharge of power.

If reflection is made, finally, on a totality of cases closed in
itself, the totality's logical concept is possibility
% --- p. 30 ---
of countless actual totalities, different in their actuality. Each
actual totality is \emph{this}, just \emph{this} totality. The
members' necessary connection, their mutual dependence and
reciprocal conditionedness, cannot be derived from the logical
concepts: connection, dependence, conditionedness. The connection is
in each totality \emph{this} connection with \emph{this} dependence,
\emph{this} conditionedness; grounded in the power-concept:
totality, the actual totality is, in other words, an antilogical
discharge of power.

Without the knowledge-concepts the power-concepts are impossible;
without the power-concepts the things are impossible, and
conversely. The infinite reciprocal determination of the logical and
the antilogical in the all-capable essence makes it then evident
from all sides that the relation between the concepts' logical
reflection in pure knowledge, their logical-antilogical reflection
in the knowing power, and the antilogical discharges in power's
immediacy is a fundamental relation originally corresponding to the
opposition of concept and existence, of possibility and actuality.

How existence's particular relations are formed and determined in
virtue of this fundamental relation, original in the essence of
power, is now to be developed, first from the standpoint of the
relative, then from that of the absolute.

\section*{b) Relative Power.}

Just as self-knowledge and knowledge of the other reciprocally
condition each other in pure ideality, so now, from the standpoint
of reality, the same must hold of self-power and power in the other.

The ideal and the real do not fall outside each other and cannot be
compounded in an outward manner. Knowledge is an ideal power, an
essential capability: to know something
% --- p. 31 ---
is, in the intellectual sense, to have power over something
(knowledge is power). Conversely, power in its immediacy is an
antilogical expression of real knowledge.

At the stage of finitude the ideal power in the form of selfhood is
a relative self-power; the object-power, the real power, is power in
the other.

\subsection*{α) Relative Self-Power.}

Since all knowledge rests on ideal capability, relative
self-knowledge is a finitely conditioned self-capability, a relative
self-power. (To swoon is to fall into powerlessness.)
% Danish: »At besvime er at falde i Afmagt« — Afmagt is both the
% ordinary word for a swoon and the technical opposite of Magt.

That the relative self-power must be conditioned by a relative power
over the other, just as self-knowledge by a corresponding knowledge
of the other, follows then directly from the unity of knowledge with
the ideal power.

Power is the essence of independence, reflected in sublated
resistance.

The resistance which knowledge's a priori representations --- such
ideal members as logic's categories, mathematics' magnitudes,
numbers, points, lines, surfaces, bodies, etc. --- offer to
consciousness affords in its immediacy so fleeting a resistance that
our knowing subjectivity can unconditionally be said to have them in
its power, so far as it can arbitrarily let them arise, alternate,
and vanish. Only when reflection is made on the represented members'
mutual relations in such wise that a system of presuppositions,
conditions, and consequences forms itself, does an ideal opposition
enter between the subjective power in consciousness's reflection and
the objective power in the members' connection. A problem is set up,
and the problem offers resistance: the thinking consciousness that
wrestles with the problem must prove its ideal power by the
sublation of the resistance, i.e.\ by the solution of the problem.

In relation to the abstractly intellectual objects there is only a
relative limit for consciousness's power, so far as
% --- p. 32 ---
there is only a relative limit for finite knowledge; but as the
human subjectivity steps up, knowing, against existence's actual
objects, it must recognize its limited power, its barrier.

The barrier shows itself in the antilogical immediacy with which the
sensible object presents itself to the sensing subjectivity. The
real is outside, the ideal inside consciousness: the relation
between inside and outside is a repelling relation. And yet the
object outside is present in its image inside: the influence that
goes out from the sensible-thinkable object becomes a moment for
that self-activity under which the sensing-thinking I sets it forth
before itself and, by this setting-forth of its other, asserts
itself. In order to see the object, the seer must be affected by the
light going out from the object; but this affecting is only a
condition that is absorbed in the self-activity under which the
image is formed. The sight-image is object-image; thus the unity
immediate in the mirroring is power-unity. And yet the seer's gaze
fastens not on the image but on the object; the unity breaks out
into opposition, and the object stands in immediate independence
before the seer.

The relative self-power is thus recognizable by the original
activity with which the knowing (intuiting, representing, thinking)
self is capable, under resistance of the other, of winning in what
is its own, yet in such wise that in this its own it recognizes its
dependence on the other. In the form of knowledge, however, the
self-power is hidden, so far as the determination of the other
everywhere has the self-determination beneath it. But as knowledge
passes over into willing, the relative self-power becomes completely
manifest.

The will's unity with the organs subordinated to it (nerves,
muscles, hand-movements, etc.) is a power-unity;
for the organic other is a moment for the determining
% --- p. 33 ---
will. Against the will's organic-actual independence the objects of
the surrounding world form an opposition in power, whose actual
unity in turn results from the self-power. The will expressed in the
self-power is in itself determined by representations and concepts;
it is in consequence capable of translating new forms, new concepts,
into the stuff of otherness (manufactures; art-productions), and in
this way reflects its independence in sublated resistance.

The fundamental difference between knowledge's theoretical and the
will's practical self-power is typically determined by the relation
between the logical and the antilogical. In theory the movement goes
from the antilogical over into the logical (from actuality to
concept); in practice, conversely, the movement goes from the
logical over into the antilogical (from concept to actuality).

\subsection*{β) Power in the Other.}

Since otherness is relativity, it follows of itself that power in
the other must be relative power. Power's distinguishing mark in
this domain is to be sought not merely in the essential difference
between concept and thing, but also in the peculiar relation between
the things themselves.

\emph{Matter} and \emph{force} are in the logical sense concepts;
but in antilogical immediacy they are actual facts. Existing matter
is selfhood's uttermost opposite, power's original discharge in the
other. The elementary parts (atoms) constitute an infinite swarm of
antilogical real-points, each with relative independence over
against the others. The atomistic members are so outwardly separated
that on superficial view it looks as though there were no connection
between them. This sundering of independent members is
characteristic of power; but to posit the independent members as
moments in active relations is, after what we have seen, also
characteristic of power: the relation between the mutually
attracting and repelling atoms is a thoroughgoing otherness-relation,
a power-relation.
% --- p. 34 ---

Relation is grounded in reflection; every relation, also the
relation between the atoms, is reflected; but in the power-relation
this reflecting is expressly to be determined as immediate activity.

Regarded for itself, the finite activity is a conditioned
manifestation of force: the reflection actualized in the
power-relation is force, a logical-antilogical category for
reflected fact. Matter for itself is power's immediate expression,
force for itself its reflected expression. The difference-unity of
matter and force is power-unity, and the power-unity in turn grounds
matter's and force's relative opposition (the force in one matter
acts on another matter; matter's selflessness: inertia and force).

\emph{Things} and \emph{laws}. The relation between the immediate
and the reflected power assumes different forms according as the
reflection of the logical in the antilogical turns. As unity of
atoms and forces the thing, the concrete-material singularity, is a
power-unity. But what the thing is in itself, it is only by what it
is in its relation to other things (that the knife is sharp means
that there is something it can cut); the unity breaks out into
opposition: opposition of thing and thing is an antilogical
reflection-opposition, an opposition of power and power, reflected
in forces and properties.

Through the particular in the reciprocity of forces the power-unity
comes into view, not in immediate and concrete form, like the thing,
but in the abstract form of affirmative reflection, as the
universal power determining the relations of things --- as the law.

That the law is power is seen not directly from its being the
universal --- for the universal in logical abstraction is as yet no
law, and the universal in logical concretion (concept, Idea) is not
either --- but from this, that the universal in the form of
universality
% --- p. 35 ---
has power over the singular, has the special cases as immediate
facts against it and yet as moments beneath it.

The connection of conditions has for its basis the thoroughgoing
power-unity, reflected in the manifold of things, resting on itself
in its generality. When the one member is posited, the other must
follow, for antecedent and consequent have their subsistence in one
and the same universal positing of power. This universal positing of
power is the thoroughgoing essentiality of the actual. Outside this
essentiality the things are nothing; they are essenceless; their
essence is the universal power, and the power is \emph{over} them.
But in and through the essentiality they are independent; their
independence is the particular power, and the power is in them. So
many systems of specially necessary member-connections, so many
forms of the determining universal, so many actual laws.

\emph{Wholes} and \emph{parts}. The thing's unity is no self-unity,
but on the contrary a unity of other with other. What the thing is
in itself, it is in its properties, and what it is in its
properties, it is in its law-determined connection with other
things: its subsistence in unity is a subsistence in plurality. The
thing-unity subsisting in a plurality of things is a wholeness, and
the things subordinated to its subsistence are parts. In the form of
the concept the relation between whole and parts is a logical
reflection-relation: when the whole is essence, the parts are
phenomena (mechanical whole), and conversely (chemical whole). But
in actuality the relation is a power-relation. The whole's
independence rests on the parts' dependence, and the parts'
independence rests in turn on a dependence in the whole. Separated
from its whole, the part is itself a whole: the whole's independence
thus rests on its dependence.

When the parts are apprehended logically as moments reflected in the
whole's unity, the existing unity --- \emph{this} whole --- must be
apprehended antilogically; conversely, when the parts ---
% --- p. 36 ---
\emph{these} existent parts --- are apprehended antilogically, the
wholeness must be apprehended logically as the parts'
connection-unity.

The unceasing converting of logical and antilogical determinations
which on the whole characterizes power in the other, and thereby all
outer, immediate actuality, has its ground in this, that the thing,
the concept's antilogical expression, is not a concept, and the
concept, the thing's logical expression, not a thing.

\subsection*{γ) The Dissolution of the Relative Power.}

The relative power has a limit over which it has no power; it has,
namely, the limit outside itself in another power. However much it
extends itself in the finite, it still cannot get the limit lifted.
Since now power's limit is its insuperable resistance, the
antilogical expression of its actual restriction, it is evident from
this that this restriction must at the same time be the logical
proof of its essential powerlessness: all relative power is
powerlessness.

In the form of reflection as well as in that of immediacy, power is
reality's expression. Immediately, the power in matter is matter's
reality, in the thing the thing's reality, just as the power
reflected in force is force's reality, etc. In logical reflection
the difference between the immediate and the reflected is a fluid
determination that sublates itself: the immediate is reflected, and
the reflected in turn immediate; but the antilogical in power fixes
the difference fast and transforms the essential transitions into
abrupt conversions.

The independent gives way everywhere and yields place to the
dependent. On what, indeed, does matter's independence rest? On
force. And force's reality? On matter. The reality is in an other,
and the reality of the other in turn in an other; the mass has no
reality in itself, its reality is in the atoms; but the atoms have
no reality in themselves either, for their reality is in the mass.
To ascribe reality to the atoms by declaring them independent parts
is of no avail, for
% --- p. 37 (the γ subsection concludes here) ---
the independent part is not a part but a whole. That the object as a
whole consists of parts is to say, in other words, that it is in the
parts' power; but the parts are in turn, by means of the connection,
in the whole's power: the whole is powerless on account of the
parts, and the parts powerless on account of the whole.

That the relative self-power is likewise powerlessness is easy to
see; it has in the other its outward limit, and in the limit its
reality. Knowledge's limit is its powerlessness; finite knowledge
cannot clarify things through and through, for the knower does not
have the actual things in his power. The will's limit is its
powerlessness; for the finite will, however strong it be, and the
strength of selfhood, however great it be, has in the outward things
a resistance it cannot overcome: its resistance is its relativity,
and the relativity its powerlessness. (Thor in Jotunheim). The
self's organically conditioned independence is its dependence; its
reality rests on the organs, but the organs' reality is
otherness-reality, and this otherness-reality is the self's
powerlessness (old age, sickness, death).

\section*{c) The Absolute Power.}

That the relative discharges of power, the finite realities,
dissolve into powerlessness at their logical-antilogical limits by
no means signifies that power vanishes, that the actual world
becomes an empty semblance; but it signifies that all finite
independence must be recognized as otherness-expression of the
infinite power. In and through the infinite power everything finite
has subsistence; as the power-unity which, in the things' relative
independence, has power over its original opposite and thereby over
itself, the infinite power is absolute.

\subsection*{α) Absolute Self-Power.}

The infinite power is not to be apprehended as the infinite
substance (Spinoza). By
% --- p. 38 ---
apprehending the infinite power as substance one makes the essence
subsisting in itself into an immediate being. But if the infinite
substance is to be the sole being, its opposite (the finite things,
the accidents) must necessarily be determined as non-being; or in
the language of reflection: if the infinite substance is to be the
sole independent essence, the existing things in their finitude must
be essenceless, consequently without any independence. And yet this
is not to be understood as though the finite things were without
being and essence; on the contrary! that which is in them is the
substance itself with its attributes: thinking and extension. Just
as the thinking beings (souls) are modifications of the attribute of
thinking, special limitations of the one thinking being, so the
extended beings (bodies) are limitations in the substance's infinite
extension.

But wherein does the essential reality consist --- where now is the
power? The infinite substance is powerless; for the unity is without
opposition. The substance is in actuality not its own cause
(\textit{causa sui}); the cause demands a substantial opposition of
cause and effect, but such a self-separation is excluded by the
substance-unity's immediate being. Over against the finite things
the infinite substance is only an empty shadow; its substantiality
is in the things; the things' limitation is its own limitation, its
own finitude: instead of being in independent unity it exists only
in infinite plurality. But the plurality is, according to the
limitation (\textit{omnis determinatio est negatio}), dependence:
far from being the substance present in all finite things, the
substance-unity has on the contrary become substanceless: the finite
negations have sucked the blood of its essence and deprived it of
all substantiality. It is not in itself, for it is immediately in
its other, and it is not in its
% --- p. 39 ---
other, for it is immediately in itself; it is neither in itself nor
in its other: it is a contradiction.

The substantial contradiction is solved as the actual doubling, the
power-doubling itself, comes into view. The infinite power-unity
consists in actual self-opposition; it is not an existent substance,
but a substantial activity, a self-activity that has its other in
itself by having it against itself: in this self-activity the
original self has infinite substantiality, and in this
substantiality consists its absolute self-power.

If the absolute self-power now annihilated the things' independence,
it would at the same time annihilate itself; it would then become a
power over shadows, a shadow-power. But if the finite substances
could subsist by themselves in outward opposition to the self-power,
then the self-power would itself be relative and, in its relativity,
powerless. Only in that the self-power, through its own opposition,
is power in the other do the existing things have relative
substantiality; their independence among themselves is grounded in
their dependence on the self-power: through the power in the other
the absolute self-power is overreaching sole power.

\subsection*{β) Sole Power and Superior Power.}

Dependence and superior power are relative concepts; but it makes an
essential difference whether the relative is determined in relation
to another relative, or is taken up and affirmed in the absolute.
Let A, B, and C be three relative members, each with its
independence. A can now gain superior power by making B and C
dependent; but this superior power, however great it be, can never
become sole power. Sole power in A would presuppose that B and C had
to such a degree lost all independence that they could offer no
resistance at all. But when all resistance vanishes, the superior
power vanishes too; when the vanquished is crushed, the mighty one
stands alone (the thrall is dependent on the master, for as existing
individual he is independent
% --- p. 40 ---
outside the master; but the dead man is not dependent on the living:
the dependence vanished with the independence). Thus: when the sole
power is posited, the superior power is sublated; when the superior
power is posited, the sole power is sublated.

Otherwise, when the relative is taken up into the absolute.

If the relative were not, neither were the absolute; if the absolute
had the relative outside itself, it were itself relative; the
relative has independence as moment in the absolute: as the absolute
is determined, its moments, the relative members, must then also be
determined. In logical reflection the absolute is the concept
completed in its comprehending self, and its relative determinations
concept-determinations, relative concepts; but in the form of power,
with the antilogical elements, the absolute is infinite
substantiality, and its real moments are relative substances.

With the relative substances' absolute dependence on the absolute
self-unity, the sole power is united with the superior power. The
superior power is sole power, for the relative substances are only
otherness-determinations of the self-power determining itself in
everything; the sole power is superior power, for the substantial
otherness-determinations are, as power's antilogical discharges,
relative substances: absolute sole power is absolute superior power.

If in this relation one-sided weight is laid on the sole power, the
absolute self-power becomes the only efficient cause (\textit{causa
efficiens}), and the relative substances --- the intellectual as
well as the material --- become the cause's substantial effects.
These effects have indeed, in finite separation, a certain
independence, so far as they are substantial; but their dependence
shows itself in this, that as effects they are not themselves
causes. The one efficient cause (the almighty will) must itself
everywhere produce the reciprocity between the effected members.
Instead of being efficient causes, the relative substances are thus
reduced to being passage-members for the all-effecting cause and
% --- p. 41 ---
thereby to being mere occasional causes (\textit{causæ
occasionales}; Malebranche).

By this conception the substances have again become substanceless;
they are again shadows: the dependence has vanished with the
independence, and the efficient sole power has sunk down into a
meaningless shadow-power.

For the sake of the dependence, weight must accordingly be laid on
the superior power. On the resistance-power presupposed in the
essence of the superior power rests the relative substances'
actuality: their actuality is that they act upon one another
mutually. Actuality consists in active relations. That two masses
exhibit attraction does not mean that the absolutely willing
self-power takes occasion from the existing members to bring them
nearer each other; it is not the attraction, but its concept, in
which the self-power expresses its superior power. It lies in the
concept of space-filling, place-taking matter that the determination
of the place one mass occupies must be reflected in the
determination of the place occupied by another mass; thus the
attraction is determined by the concept. But the concept does not
act as concept; the concept of attraction is not attracting; the
thing is cause, but the cause is not thing. It is the masses which,
in virtue of the cause-concept, attract each other reciprocally; but
the cause-concept does not attach to the thing as if the thing could
be thing without being cause; the existing thing is the concept
stepping out into power, converted into existence: the existing
thing is efficient cause.

The things are absolutely dependent on the concepts; in this
dependence they have their independence; but the concepts are,
according to their origin, absolutely dependent on the comprehender:
the comprehending sole power thus has, in virtue of the concepts,
infinite superior power over the thing.

\subsection*{γ) The Idea of Power.}

The concept is the thing's ideal, and
% --- p. 42 ---
the thing the concept's real, power. The unity of the ideal and the
real power is the Idea of Power.

Apprehended for itself in logical generality, the power-concept is a
reflection-category, a peculiar knowledge-form, a one-sided concept
among several others. But, seen in the Idea, power's infinity
receives its light from knowledge. Every power-concept is a
knowledge-concept, but not conversely. The knowledge-concept for
itself is powerless; it reflects the thing in the thing's semblance;
its reflecting is only a shining, a mirroring. The power-concept, on
the other hand, whose reflecting of antilogical moments is
positively posited as activity, is actuality-capacity, capacity for
actuality, the actual's real possibility. The sharply marked
opposition of concept and thing, of possibility and actuality, is,
as we have seen, an opposition within power in power itself, an
opposition of negatively reflected and positively immediate power.

The real-concept then shows itself, from the Idea's standpoint, in a
double-sided relation to the thing; it is indwelling in the thing
and yet infinitely overreaching. The concept of sulphur, for
example, is indwelling in sulphur; its existence is the sulphur
existing in a given state. And yet sulphur in antilogical-immediate
existence is infinitely far from expressing the whole concept. To
exhaust the content of its concept, sulphur must pass through all
possible states, show itself in all possible combinations; but
actuality permits each time only \emph{this} state, \emph{this}
determinate combination: actuality, the antilogical determinateness,
is finite; but possibility, the logical essentiality, is infinite.
Sulphur's concept is not merely reflected in a peculiar circle of
kindred stuff-concepts; it is, through the stuff-concepts'
reflection, originally reflected in the whole infinite system of
concept-systems: as Idea-form the finite concept is infinite.

The fundamental relation between the indwelling and the overreaching
real-concept unfolds itself in an opposition of two heterogeneous
% --- p. 43 ---
circles of movement, the ideal and the real, whose unity in the Idea
is a thoroughgoing inference-unity.

In the ideal circle the movement goes through real-negations of
antilogical reflexes from concept to concept, without the ideal
activity at any point coming to discharge in outer actuality. Such
stuff-concepts as oxygen, hydrogen, carbon, gold, etc., can be
comprehended only by being reflected in a system of logical
universal concepts. What such a reflection has to signify can be
seen from the immanent logic's dialectical analyses. The
stuff-concepts named are special modifications of the
substance-concept. The special arises by regard to the stuffs
themselves, hence by reflexes of the antilogical; but the emphasis
falls not on the existing substances' peculiarity, but on their
universal essentiality, i.e.\ on \emph{the concept: substance}. The
substance-concept is the logically developed essence-concept, the
essence-concept the logically developed being-concept. The movement
goes synthetically: from being through essence to substance;
analytically: from substance through essence to being.

That the development of the categories which the immanent logic is
able to give is by no means the absolute knowledge's own unfolding,
but only a formal-dialectical dissolution of the finite reflection
in the infinite, has already been shown. Nonetheless there is in the
development a reflected shine of the eternal logic, a reflected
splendour of the divine reason, whose inexhaustible depth is
selfhood.

In the real circle the movement takes place from actuality through
possibility to actuality, in that the possibility everywhere rests
on the indwelling concept's logical overreaching beyond the finite
limit, while actuality indicates the limit antilogically. Actuality
of given things is possibility of other things (the actuality of
oxygen and hydrogen is the possibility of water); the by-things
co-operating with the principal things afford conditions.

% --- p. 44 ---
Through possibility the things by no means go beyond their own
concepts, but they overstep the actuality-limit to which the
existent's antilogical determinateness is bound. If there were not
more in the thing's concept than there immediately is in the
existing thing, the thing could not perish; if there were not more
in the actual's concept than in the actual itself, the actuality of
the one could not be possibility of the other; but if the one's
actuality were not the other's possibility, nothing could be
grounded: no thing could arise out of the ground.

The fact that the things' possibilities are hidden under the forms
of actuality has occasioned the misunderstanding that possibility
did not actually exist (the Megarics' assertion), a misunderstanding
that only proves the impossibility of making things into concepts
and concepts into things. It is precisely \emph{the secret of the
Idea of Power} that in the world of things there are no concepts,
because the things themselves, in their being, becoming, and
transitions, exhaust the concepts.

The two circles now form two worlds, a world of concepts, \emph{the
ideal system}, and a world of things, \emph{the real system}. From
the one of these two worlds there is no finite transition to the
other: from concepts come concepts, from things things, into the
infinite. The transition from concept to existence cannot possibly
be thought in such wise that the all-comprehending self should first
specially think out for itself a special concept and thereupon apply
power in order to produce the object corresponding to the concept.
If the matter stood thus, power would be arbitrarily joined on to
knowledge, the things' actual connection would be broken, and the
real possibility, conditioned by the connection, groundless (the
rose, for example, would then by a fiat of power step forth
immediately out of the rose's concept, the horse
% --- p. 45 ---
out of the horse's concept, etc.); but a fiat of power that breaks
the things' connection sublates all logic.

The fundamental concept that expresses the concept's and actuality's
original unity in power is the concept \emph{disposition}.
% Danish: Anlæg (germ, endowment, anlage) — rendered »disposition«
% throughout this section.

In the world of things the reciprocal determination of actuality and
possibility runs out into endless series: the actuality of A is the
possibility of B; as A goes to the ground, B comes forth out of the
ground, steps out into existence and becomes an actuality; but the
actuality of B is in turn the possibility of C, and so forth. What
now is herein the originally determining? It cannot be the actual,
for the actuality in every succeeding member is precisely determined
by the possibility lying in the preceding one. Nor can it be the
possible, for the possibility in every preceding member is expressly
determined by its own actuality. The fundamentally determining must
necessarily be sought in a concept which in the form of possibility
has taken up all the moments of actuality, and thus in and through
the possibility has predetermined actuality. But the predetermination
is possible only in that the concept in the form of actuality has
predetermined all becoming possibilities. On account of the
possibilities' predetermined unity, the concept has power over the
actual's manifold. The unity of the possibilities is at bottom
itself the actual's own presupposed unity; with actuality's unity as
basis the concept has power over the possibilities' manifold: this
power-concept, characteristic by its doubling of positing and
presupposing, is the disposition.

In the disposition power is concept, and the concept power. To
reflect on the concept is to determine the disposition's relation to
possibility; to reflect on the power is to determine its relation to
actuality.

\emph{Disposition and possibility.} Every disposition is a
possibility, but not every possibility is a disposition. The
difference rests on the conditions. In the outward, finite world
% --- p. 46 ---
the relations are outward, and from the outward relations follow
\emph{outward} conditions. Now so far as the real possibilities
indwelling in the things depend on outer conditions, they are
possibilities pure and simple, not dispositions (possibility of a
rainbow; possibility of rain and wind, etc.). In the essence of
things, in the system of real-concepts, the relations are on the
contrary inward; to the inward demands lying in the determinate
real-concepts correspond the \emph{inward} conditions: the
possibility originally determined by its inward conditions is a
disposition (the germ in the seed).

\emph{Disposition and actuality.} Every disposition is an actuality
(the seed; the animal semen); but not conversely. To the possibility
with the outward conditions corresponds the outward actual in which
there is no disposition (a stone barrow); the actual in which there
is disposition is a selfhood-product, is expressly to be determined
by indwelling selfhood (organ, organism). To be sure, the merely
outward, selfless actual can be apprehended from the standpoint of
selfhood, and its indwelling possibility so far determined as
disposition. Thus one can say, for example, with respect to water,
that in oxygen there is a disposition to combination with hydrogen,
and in hydrogen a disposition to combination with oxygen. But this
by no means proves that the selfless outward is in itself a
disposition; it shows only that the finitude in the outward actual
is posited by the self as passage-member for the overreaching
concept, and thereby reduced to a condition for satisfying a demand
of selfhood upon the other.

\emph{Disposition and end.} The absolute concept-demand is a
selfhood-demand, a totality-demand closed in itself; the demand
expressly presupposes an other in which and whereby it is to be
fulfilled: between the demand and the fulfilment lies the
disposition.

The disposition is the fulfilment's real possibility, the demand's
first, immediate satisfaction. But this immediacy
% --- p. 47 ---
is in itself unsatisfying; for it conflicts with the disposition's
essence to remain in the immediate state of disposition. The first,
original disposition expresses, namely, not merely the immediate
actuality's dependence on the determining totality-concept; it
expresses at the same time that the totality, whose actuality is as
yet only a possibility, is to be developed to full actuality: the
fully developed totality is the demand's last, absolute
satisfaction, a satisfaction that can be reached only by the
roundabout way of mediacy. Hereby the concept disposition passes
over, from the subjective as well as from the objective side, into
the concept \emph{end}.

Subjectively the end is posited in the totality-concept
corresponding to the selfhood-demand. According to its free,
overreaching form this concept, determined by the comprehending
selfhood, belongs to the ideal system; but in virtue of the
disposition it must, as we have seen, belong just as originally to
the real system, for the disposition's possibility presupposes
actuality. Through the disposition the transition takes place from
the ideal to the real system, hence not in a finite but in an
infinite manner: it takes place through the end (first extreme).

In order to determine the whole actual totality, the concept must
beforehand be determined as end. In and through the end-concept the
real-concepts are prior to the things (\textit{universalia ante
rem}); but in virtue of the actual end they are at the same time
indwelling in the things (\textit{universalia in re}). The
totality's development out of the disposition is the parts', the
particular members', development: the finite comes forth out of the
infinite; thus arise the means (second extreme).

The means in their immediate singularity are antilogical; the end in
its universal essentiality is logical. But the end is actual only as
the means' end: the logical in the end is thus, in virtue of the
means, antilogical. Conversely, the world of things with its
possibilities, conditions, and actuality-sides is a world of means;
these
% --- p. 48 ---
means are the end's own means: the antilogical in the means is thus,
in virtue of the end, logical (the inference's unity is unity in
power).

In the real world, the actual totality of end and means, the
reciprocal determination of the logical and the antilogical in the
essence of power is infinitely articulated: the Idea of Power is the
articulated form-determinations' absolute content.

% Part B ends here on p. 48; beneath the two closing paragraphs a
% short centered rule, the rest of the page blank (carried as a
% comment here, as in the transcription).

\chapter*{C. Truth.}
\markboth{C. Truth.}{C. Truth.}
% --- p. 49 ---
The essential difference between knowledge and power develops itself
through the ideal and the real system into complete opposition; but
the Idea is in both systems the unity of its own opposition, and by
this unity is everywhere to be determined as truth.

Everything, even the untrue, has its truth (a true lie); and yet it
is precisely truth itself that makes untruth manifest
(\textit{verum est index sui et falsi}).

In the form of reflection, truth is unity of thinking and being:
thinking has its truth in being, and being in thinking; untruth is
accordingly discord between thinking and being. In the form of
judgment, truth is identity of subject and predicate: the object
(the subject) has its truth in the concept (the fully developed
predicate), and conversely; untruth, the contradiction in the
judgment, is accordingly strife between subject and predicate. In
the form of inference, truth is the middle of the extremes: the more
independent the extremes, the fuller the middle, and the deeper the
truth; untruth, on the other hand, is the middle's dissolution into
the mutually devouring extremes (spiritualism --- materialism).

The logical development of the Idea of Truth's peculiar fundamental
forms has in selfhood the principle of subjective truth, in
otherness the principle of objective truth,
% --- p. 50 ---
and in the self-unity infinitely returning through its
otherness-totality the principle of absolute truth.

The threefold truth-principle has, in the subjective as well as in
the objective extreme, laws of development for the whole truth
(dualism); but the one-sidedness with which each of the extremes is
burdened must be lifted in the original middle in which they come to
agreement. The extremes are not equalized by being dissolved, so
that the middle becomes merely the neutral; their equalization is
that the opposition is preserved in the unity (identity-dualism).
The inference-unity in which the opposition subsists is then itself
the absolute beginning-unity from which it proceeds.

\section*{a) Subjective Truth.}

The principle is selfhood; hence the logical emphasis falls on the
independence which subjectivity is in a position to found on
otherness's overcome resistance, in that, in overcoming the
resistance, it grounds itself. The self-grounding shows itself at
different depths as we, beginning with the finite subjectivity, work
our way forward, through a dialectical dissolution of finitude's
barriers, to the concept of the absolute subjectivity.

\subsection*{α) Finite Subjectivity.}

The judging self is the measure of things; by this proposition an
essential mark of subjective truth is indicated. But the proposition
becomes untrue when the judgment is to concern the essence of things
and rest on the immediate subject's contingent mood (Protagoras).
The proposition's truth depends on subjectivity's truth;
subjectivity's truth depends on its judgment about itself; but its
judgment about itself depends in turn on the force of inwardness
with which, in original self-positing, it has been able to ground
the independence of its essence.

Subjectivity is no immediately given object,
% --- p. 51 ---
no object subsisting for itself; the object arising through
self-objectivation is the subject itself: the truth in the object
rests on the truth in the subject, hence on the selfhood-energy with
which the objectivation goes on. That the selfhood-energy is the
truth means that subjectivity itself makes itself into the truth ---
that, in other words, truth is unconditionally subjective.

The finite subjectivity's contradiction now shows itself in this,
that as subjectivity it is to posit the whole subjective truth in
itself, and yet as finite subjectivity, restricted by finite objects
and objective conditions, it must have the truth in an other than
itself: its activity and striving must accordingly aim at solving
this contradiction.

Subjectivity's finitude is seen in its strife with the barriers. The
finitude rests not merely on the self's restriction by the outer
objective; also in its inner, in the relation between its thinking
and willing, the finitude is seen. There is a willing in all
thinking, and a thinking in all willing; but the opposition of
thought and will steps forth in all its finitude when subjectivity
is to posit the truth in its own essence; for subjectivity's truth,
the most subjective element in subjectivity's self-positing, is its
will.

In order that the truth may be posited in the will, the will must
lay thinking under requisition: the categories run on deliberation
and choice. In deliberation, thinking labours not to see into the
validity of things in themselves, but only to find out the
possibilities decisive for the will. When from the circumstances at
hand it has found out two or more possible cases, the will is to
choose. How deep-going the difference between subjective and
objective thinking is will in this connection easily be seen.

Faced with an objective theorem, thinking can indeed begin
subjectively, so far as it can begin by deliberating how the proof
may best be conducted; but that the conduct
% --- p. 52 ---
of the proof is itself objective is seen from this, that the proof
ends by making choice impossible. The proven theorem must be
accepted as it is; here there is no choice at all: the necessity is
objective. It is the converse with deliberation; it does not end
with an irrefutable proof that one of the several cases possible for
the will is the only necessary one. On the contrary! The more
decisive the question of will is, the plainer it becomes that for
each of the alternatives set up there are grounds both against and
for. The grounds for can be outweighed into the infinite by the
grounds against, and conversely. To end the deliberation, endless in
reflection, the will must choose. The choice is necessary, and when
the crisis is great, unconditionally necessary; but the necessity is
subjective.

Whether now the will under the given circumstances has made a true
or an untrue choice cannot be decided objectively by regard to the
final outcome, for the outcome is, in the moment of deliberation,
only a possibility, not an actuality. What the outcome can at length
clear up is, in the objective sense, solely whether the choice has
been fortunate or unfortunate, beneficial or harmful for the matter
concerned; but these finitude-determinations still decide nothing
about the truth in the choosing subjectivity. \emph{To choose is to
venture}: the truth rests on the depth of self-understanding in
which the chooser has recognized his \emph{responsibility}.

To take responsibility upon oneself is to posit the truth in the
self; the strength of inwardness with which this happens receives
its expression in the resolution with which the choice is held fast.

Subjective truth is related to objective truth as resolution to
inference.
% Danish wordplay: Beslutning (resolution) / Slutning (inference);
% below, »Afslutningen« (the closing) continues the same root.
The conclusion (\textit{conclusum}) drawn from an inference's
presuppositions (premises) may not contain more than what is
incontestably grounded in the presuppositions: this is the
conclusion's objective necessity; in this necessity it has its
truth.
% --- p. 53 ---
Of the resolution the opposite holds. As the resolution is taken,
the will's self-positing adds something of its own, something not to
be found in the deliberation's premises, and this something, the
closing itself, is precisely the decisive element.

In the objective inference it is the premises that must vouch for
the conclusion's validity; therefore the rational act is ended with
the inference itself. In the resolution, the subjective inference,
the will's self-positing is on the contrary a positing of the
consequent which, under the weight of responsibility, must vouch for
the premises' validity. The resolution, which is borne only by
grounds it must itself bear, is at its maximum an antirational act.
The antirational in the resolution's inner actuality corresponds
exactly to the antilogical in the thing's outer actuality: in the
friction between the inner and the outer actuality the resolution
steps out, to be objectivated in actual action.

How the finite subjectivity, in the endeavour to posit the truth in
itself, must step by step struggle with contradictions is then
obvious. To make a choice, it must lean on deliberation; but
deliberation's contradiction is that its grounds for and against are
not grounding. Only the choice is capable of making a ground
insufficient in itself sufficient as motive: a motive is a
subjective ground, a ground the will has adopted on its own
responsibility. The contradiction that the will both has and yet
does not have a truly grounding ground the resolution must strike
down as action is proceeded to. But the contradiction returns with
redoubled force. The subject in its finitude, namely, does not take
the resolution unconditionally for its own sake, but also for the
action's sake. The action's outcome is only in part dependent on the
agent, and yet the agent has taken the responsibility upon himself.
The outward-going action, which was to prove the truth of the choice
and therewith of the resolution, can then precisely by its outcome
turn
% --- p. 54 ---
against the acting subject and convince it that the truth
represented in choice and resolution was after all an untruth.

\subsection*{β) Subjectivity's Abstract Infinity.}

The subject's finitude is its untruth; if this proposition is held
fast by itself, the task must become to infinitize subjectivity. As
existing individual the subject is subjected to finitude's terms; it
lacks the power to subject the outer concrete actuality
unconditionally to itself. The infinitizing must accordingly begin
with the self, by a peculiar practical abstraction, freeing itself
from the outward and contingent and declaring everything that does
not stand in its power indifferent; the truth then lies not in
things and events, but in the self's universal essence, in the
thinking-willing subject's agreement with itself.

The existing subject has thus placed itself at the standpoint of
infinity; deliberation is transformed into essential
self-recollection, and choice into unconditional self-choice.

In the self-recollection subjectivity then does not begin by
clearing up its position toward these or those actuality-tasks,
finding out these or those alternatives --- all such concerns only
finitude, and finitude is untruth --- no, it will in the infinity of
its thinking itself be the truth; so the subject raises itself above
its individual limitation, converts its natural singularity into a
thought singularity, and leans on the logical judgment: \emph{I, the
singular, am the universal}.

To the infinity of thinking corresponds the infinity of the will.
That an outward-going will wills in the world of things this or that
singular is of subordinate significance. Just as the will in its
natural singularity is untrue, so also everything outwardly singular
at which the will aims is something in itself indifferent and
untrue. \emph{The true singular will is the universal will}; the
universal will's object is not the
% --- p. 55 ---
outward but the inward singular, is the universal, willing self, for
only in the self as thinking is the singular itself the universal.

In its universal, liberated self-being, subjectivity is now to set
about sublating the action's finitude and transforming the
resolution into an adequate inference. And this can indeed very well
be done --- namely, by the help of abstraction.

Thinking, which at this stage is no longer practically finite but
theoretically infinite, is a thinking of the thinking essence's
universal validity; the content of selfhood is clarified in the
light of truth: subjectivity becomes transparent to itself. The
transparency is conditioned by the thought-essence originally
positing its universal being in the form of thinking; this is the
self-energy's first extreme.

The will, which at abstraction's high point has ceased to be a
finite, arbitrarily choosing singular will, is, though one with
thinking, yet different from thinking. Thinking strives from the
singular toward the universal, the will conversely from the
universal toward the singular: it is the same striving, and yet a
different striving. On account of the difference between thinking
and will, the content of selfhood must then pass over from the form
of thinking into the will, must by a peculiar self-energy be
determined as the will's content; thus the second extreme forms
itself.

The middle of the extremes is the ground itself from which they
proceed; it is the thinking-willing subjectivity's original
self-unity; in this unity's self-development they have their truth.
The ground-division corresponding to the judgment-act yields two
judgments: \emph{the self's universal content is a thought-content},
and \emph{the self's universal content is a will-content}; but the
two judgments are premises for one inference: the inference-act is
the self-activity gathered in its universal, the pure logical act,
freed from finitude.
% --- p. 56 ---
In the closing-together of the extremes the resolution is then
transformed into an inference necessary in itself, logically
adequate. The necessary form-doubling whereby the extremes arise is
grounded in the essence-unity in which they go together as into
their original middle. There lies then nothing at all in the
concluding proposition beyond what is contained in the premises; at
the transition from the form of thinking to the form of will, the
will has nothing to add: it has no choice, no risk, no
responsibility, for the transition takes place with necessity. The
will wills what it \emph{must} will: freedom is necessity; but what
the will must will, it wills of itself from its essence's inner
prompting: necessity is freedom (\textit{libera necessitas}).

But just as subjectivity, in this transparent self-being of its
essence, has gathered itself out of the dim manifold of life's
distractions and thrust the finite from itself, it gets in earnest
to do with finitude.

The transparent self-being is an infinitely empty being; every
moment of difference, every condition of concretion, every least
gleam of content it owes to reflexes of other-being. The whole
liberation, founded on necessity, is conditioned by an energetic
abstraction; but the abstraction's energy is precisely conditioned
by the finitude from which abstraction is made. If subjectivity will
hold the abstraction fast in order to keep finitude out, it becomes
itself a hollow abstractum, a metaphysical shadow-essence. If for
concretion's sake it will engage with finitude, then the cleft
between the nature-bound singular, the existing individual, and the
singular of thought, the selfhood transparent to itself, becomes so
wide that the whole of finitude streams in, fills and overwhelms the
abstract-universal subjectivity.

The individual's universal will cannot possibly come to existence
except through its determinate singular will. The object of this
singular will is not the universal itself, but a
% --- p. 57 ---
singularity existing in the world of things, in which the universal
is to be expressed. (No one can in all generality do the general
duty; but the singular person can and ought in each single case to
do his duty). So practical thinking parts again from theoretical; so
deliberation, choice, resolution, with all the contradictions
pointed out in them, begin over again, and subjectivity is no longer
transparent to itself.

The universal willing of the universal is unactual, and the same
holds of the pure, universal acting. The actual action is, even when
it belongs to inwardness, invariably an action in the particular by
means of the given singular, hence a particular action. The truth in
the particular action now depends on: 1) that the action in its
singularity is the will's complete expression; 2) that in its
connection with other particular actions, through which the acting
subjectivity, in solving a determinate actuality-task, will
unconditionally satisfy itself, it corresponds exactly to the aim.
If these two conditions are to be fulfilled, subjectivity must have
the particular action, its actual execution, its outcome and
consequences, in its power; here no abstracting helps, no creeping
into shelter behind the universal: subjectivity has the
responsibility.

To declare the finitude-elements that do not stand in the
individual's power contingent, unessential, and indifferent is an
empty evasion. If subjectivity cannot get its actual will actually
carried through: then the concrete actuality has broken through the
singular will's universal, and truth demands a confession that the
alleged self-agreement is an actual self-contradiction. If
subjectivity will extricate itself from the entanglements by letting
the actual tasks drop and beginning over with a willing of the
universal in the universal, it falls back into the abstraction, and
% --- p. 58 ---
truth requires a confession that the alleged self-agreement is an
essential contradiction.

If the universal in all actual actions, the pure, universal acting,
is to be the only essential, then the actual actions in their
singularity become unessential; and the truth, or rather the
untruth, is that the universal action itself becomes unactual on
account of the abstraction. If conversely the actual actions, the
universal action's actual expression, are made the essential on
account of the concretion: then the finite gets power over the
infinite; the circumstances in the singular overwhelm the will in
the universal, and the untrue subjectivity, entangled in
self-contradictions, collapses in powerlessness.

\subsection*{γ) The Absolute Subjectivity.}

Subjectivity's untruth has shown itself from two principal sides:
first, in its finitude it will posit itself as the truth --- that is
a glaring contradiction; next, when the self-contradiction has
become manifest, it will by a radical abstraction, despite its
finitude, make itself one with the infinite: the latter is in the
logical sense worse than the former. The consequence is that the
immanent logic reduces the subjective to a moment sublated in the
Idea and deifies objectivity. The Idea is indeed to be the unity of
the subjective and the objective; but the universal is the
objective: the overreaching universal must then make truth itself
absolutely objective.

The fundamental contradiction is obvious. If the original
subjectivity is not itself absolute, there is no absolutely knowing
self; if there is no absolutely knowing self, there is no absolute
knowledge; if there is no absolute knowledge, there is no absolute
truth. Therefore: in the name of truth the eternal, unoriginated
subjectivity must be absolute.

Hereby objectivity is neither excluded nor weakened; objectivity is
in the power; the absolute subjectivity is knowledge's and power's
original self-unity: the absolute
% --- p. 59 ---
subjectivity is the truth. But self-unity of knowledge and power is
will-unity; the task becomes accordingly: to determine the concept
of the absolute will.

How does the relative proceed from the absolute? With absolute
necessity! This answer is \textit{correct} so long as the
fundamental relation is apprehended in all generality, for as the
relative's possibility the absolute is itself impossible without the
relative. But when the question \textit{in concreto} is about this
or that antilogically determined relative, the relation becomes
quite another. In its proceeding from the absolute, namely, the
relative is itself a relative absolute, a totality in a peculiar
disposition. From the disposition-concept in general to the
disposition determined in actuality there is a leap that no
necessity-concept can fill.

If the absolute in itself were nothing other than the dim ground
overfilled with existence-content, then the absolute necessity with
which the existent broke forth out of the ground would be one with
blind contingency. But to begin with blind contingency is to begin
as illogically as possible. Therefore: how does the general
totality-concept pass over into becoming this, just this
power-concept, corresponding to this, just this totality --- in
other words: how does the disposition, determinable in all
indeterminacy, pass over into being determined? To this decisive
question logic has only one decisive answer: it happens in virtue of
an absolute choice.

The finite subjectivity has only a choice between two or more
possibilities restricted by finitude's conditions; deliberation must
go before, to find out what under the given circumstances may be
regarded as the best: the choice between a relatively better and
worse is a relative choice. The absolute subjectivity, on the other
hand, has, as regards the disposition, not merely two or three but
countless possibilities; for otherness, the infinite stuff of
actuality, receives
% --- p. 60 ---
all its determinateness from the determining real-concept; but the
real-concept itself must originally be determined by the
all-comprehending self-determining: for the unconditioned
self-determining, the absolute will, everything is possible.

Subjectivity's absolute superiority over the disposition-task shows
itself just in this, that the task can be solved in countless ways;
the one of the alternating possibilities is, for the infinite
comprehending, not better or worse, not dimmer or clearer than the
other: they are in the depth of grounds equally grounded, equally
clear, equally possible. But if all the disjunctive judgment's
alternating possibilities are equally possible, then, did the choice
not enter, they would all be equally impossible: in order not to
drown in a sea of possibilities, the absolute subjectivity must
absolutely choose this or that determinate possibility and express
the disposition-concept in this or that determinate disposition.

The choice is necessary in the absolute in so far as it is
absolutely necessary that a choice be made; but what is chosen ---
which of the possible cases, the possibilities equally grounded in
the concept, is chosen --- rests on the will alone. When the
absolute subjectivity then chooses this, just this possibility,
this, just this disposition, the will satisfies itself
unconditionally. This is the rational-antirational that makes the
self-determination into will; it is the divine, absolutely thorough
arbitrariness: \textit{sic volo, sic jubeo, stat pro ratione
voluntas}.

The absolute choice is one with the absolute resolution. At the
stage of finitude it looks like a defect in the resolution that it
cannot be built on objective assurance, cannot coincide with an
inference grounded in itself. From the absolute standpoint the
relation is reversed. It is precisely a perfection in the resolution
that it takes the inference up into its ground and stands above the
inference. The inference belongs to the concept as concept, and for
that
% --- p. 61 ---
reason can never reach beyond the universal essentiality of the
concept-systems. That the unity of disposition and execution, of end
and means, is in actuality an inference-unity is to say, in other
words, that it is a logical-antilogical concept-unity.

But the antilogical in the outer actuality corresponds, as already
shown, to the antirational in subjectivity's inner actuality, to the
will's choice --- hence to that which completes the inference in
resolution.

If the resolution could be torn loose from the inference, then the
absolute subjectivity could at the same time tear itself loose from
the absolute reason, and the antirational, the counter-rational,
would be what conflicts with reason. But the resolution expresses
just the converse; it expresses, namely, that subjectivity precisely
satisfies reason's demand in that, by applying the factor of
self-determination, it integrates the universal reason-grounds and
transforms them into will-grounds\footnote{How the opposition
between two absolutely heterogeneous principles of cognition
proceeds from this, cf. the author's Religionsphilosophie Side
28---35; 76--81, and several other places.}.
% The Danish footnote prints »jvnf. Forftrs. Religionsphilosophie«
% (Forftrs. = Forfatterens, the author's); page ranges as printed
% there (em-dash, then double hyphen).

The absolute resolution is one with the absolute action; by means of
the absolute action, having power over the conditions, the will in
the absolute is the original cause of everything relative.

This cause of causes is in the action itself at once both ideal and
real. It is \emph{ideal}; its ideality is its merging in the system
of grounds, its agreement with the eternal reason in the absolute
knowledge. It is \emph{real}; its reality is the self-power's unity
with the power in the other. The will in the absolute action is the
all-articulating actuality-principle, the principle that in the
world of things makes the ground into cause and the cause into
% --- p. 62 (the γ subsection concludes here) ---
ground; thus the will is the all-capable objectivity-principle.

The absolute action is the complete expression of the will's
activity in the realm of power. It utters itself in the relative, in
everything finite; but it does not begin finitely here or there. In
virtue of its unconditionally subjective essence, it aims decisively
at the objective. It is the universal acting, for it aims at the
universal: it is in the end the overreaching that determines the
means, in the real-concepts the energy that determines the things.
It is an infinite system of particular actions, for it aims at the
particular: it is the genus's energy in the species, the species' in
the individuals. It aims at the singular, for it confirms
otherness's right, the relative substances' antilogical
independence: it makes the things into causes, the parts into means,
the individuals into bearers of the species' disposition.

In the absolute subjectivity the truth is absolute, but absolutely
subjective.

% [translation continues from p. 62, at \section*{b) Objective Truth}]
% [text to be added: pp. 62--76]
% [text to be added: pp. 77--89]

\backmatter

\chapter*{Contents}
\markboth{Contents}{Contents}
% [text to be added: fol. XIII--XIV]

\end{document}
